\documentclass[11pt,a4paper]{article}
\usepackage[utf8]{inputenc}
\usepackage[czech]{babel}
\usepackage[left=2cm,right=2cm,top=2.5cm,bottom=2.5cm]{geometry}
\usepackage[dvipsnames,x11names,svgnames]{xcolor} 
\usepackage{graphicx}
\usepackage[hidelinks]{hyperref}
\usepackage{tocloft}
\renewcommand{\cftsecfont}{\small\sffamily}
\renewcommand{\cftsecpagefont}{\small\sffamily}
\setlength{\cftbeforesecskip}{1pt}
\usepackage{recipebox}

\begin{document}

% ==========================================
% TITULNÍ STRANA (ULTRA STABILNÍ VERZE)
% ==========================================
\begin{titlepage}
    \begin{tcolorbox}[
        enhanced,
        height=\textheight,
        colframe=Mahogany,
        colback=Orange!5,
        arc=12mm,
        boxrule=4pt,
        halign=center,
        valign=center,
        % Pokud by shadow stále zlobil, stačí tento řádek smazat
        drop shadow={black!30} 
    ]
    \vspace*{2cm}
    {\Huge \bfseries \color{Mahogany} POSBÍRANÉ RECEPTY \par}
    
    \vspace{1cm}
    {\Large \itshape \color{OliveGreen} pro Řinečku a Zaníka \par}
    
    \vspace{3cm}
    % \includegraphics[width=0.7\textwidth]{obrazek.jpg}
    \vspace{3cm}
    
    \vfill
    {\Large \bfseries \color{Mahogany} Rodinná Kuchařka \par}
    \vspace*{1cm}
    \end{tcolorbox}
\end{titlepage}

% ==========================================
% OBSAH A ZKRATKY
% ==========================================
\pagenumbering{Alph} % Dočasné číslování pro titulku (opraví varování logu)
\tableofcontents
\newpage
\pagenumbering{arabic} % Ostré číslování receptů od strany 1

\section*{Zkratky:}
\begin{itemize}[label=\color{Mahogany}\faChevronRight]
    \item \textbf{PL} = polévková lžíce
    \item \textbf{KL} = kávová lžička
\end{itemize}
\vspace{1cm}

% --- Tady už pokračují tvoje recepty ---
% ==========================================
% RECEPTY – POLÉVKY
% ==========================================

\begin{recipe}[NavyBlue]{Jíška}
    \begin{ingredients}
        \item \parbox[t]{\linewidth}{\raggedright \piece \textbf{1 cm}\\ Másla}
        \item \parbox[t]{\linewidth}{\raggedright \spoon \textbf{1,5 PL}\\ Hladké mouky}
        \item \parbox[t]{\linewidth}{\raggedright \liquid \textbf{dle potřeby}\\ Voda nebo mléko}
    \end{ingredients}
    \begin{steps}
        \item Do rendlíku dám máslo a nechám na plotně pomalu rozpustit. 
        \item Když je máslo rozpuštěné, přidám hladkou mouku a míchám na mírném ohni. 
        \item Když je jíška zlatavá a krásně voní, přidám trošku studené vody nebo mléka a zase míchám. Když se spojí voda s moukou, zase přidám trošku vody a míchám. 
        \item Tento proces opakuji tak dlouho, až má jíška konzistenci řidší kaše. Jíška nesmí mít hrudky, pokud se tak stane, přecedím jí přes jemné sítko.
    \end{steps}
\end{recipe}

\begin{recipe}[ForestGreen]{Polévka písmenková}
    \begin{ingredients}
        \item \parbox[t]{\linewidth}{\raggedright \liquid \textbf{2 l}\\ Vývaru (hovězí/vepřový/kuřecí)}
        \item \parbox[t]{\linewidth}{\raggedright \indefinite \textbf{2 hrsti}\\ Zeleniny}
        \item \parbox[t]{\linewidth}{\raggedright \piece \textbf{1 hrnek}\\ Písmenek (těstoviny)}
        \item \parbox[t]{\linewidth}{\raggedright \spoon \textbf{1 PL}\\ Bylin (petržel, libeček)}
        \item \parbox[t]{\linewidth}{\raggedright \weight \textbf{dle chuti}\\ Maso z vývaru}
        \item \parbox[t]{\linewidth}{\raggedright \piece \textbf{1 stroužek}\\ Česneku}
        \item \parbox[t]{\linewidth}{\raggedright \indefinite \textbf{špetka}\\ Pepře}
    \end{ingredients}
    \begin{steps}
        \item Dáme vařit 1 – 1,5 litru vody, když se vaří, osolíme a nasypeme písmenka. Uvaříme je podle návodu. 
        \item Když jsou písmenka měkká, scedíme je a pod studenou vodou propláchneme. Dáme stranou okapat. 
        \item V jiném hrnci přivedeme k varu vývar. Přidáme zeleninu a maso a vaříme do měkka (5-10 minut). 
        \item Přidáme okapaná písmenka a převaříme. Odstavíme z ohně, opepříme, přidáme rozdrcený česnek a bylinky.
    \end{steps}
\end{recipe}

\begin{recipe}[ForestGreen]{Polévka s nudličkami}
    \begin{ingredients}
        \item \parbox[t]{\linewidth}{\raggedright \liquid \textbf{2 l}\\ Vývaru}
        \item \parbox[t]{\linewidth}{\raggedright \indefinite \textbf{2 hrsti}\\ Zeleniny}
        \item \parbox[t]{\linewidth}{\raggedright \piece \textbf{1 hrnek}\\ Nudliček}
        \item \parbox[t]{\linewidth}{\raggedright \spoon \textbf{1 PL}\\ Bylin}
    \end{ingredients}
    \begin{steps}
        \item Postup je stejný jako u polévky písmenkové, jen zaměníme písmenka za nudličky.
    \end{steps}
\end{recipe}

\begin{recipe}[ForestGreen]{Polévka s kapáním}
    \begin{ingredients}
        \item \parbox[t]{\linewidth}{\raggedright \piece \textbf{2 ks}\\ Vejce}
        \item \parbox[t]{\linewidth}{\raggedright \spoon \textbf{1 PL}\\ Mléka}
        \item \parbox[t]{\linewidth}{\raggedright \indefinite \textbf{dle chuti}\\ Sůl, pepř, bylinky}
        \item \parbox[t]{\linewidth}{\raggedright \weight \textbf{dle potřeby}\\ Polohrubá mouka}
    \end{ingredients}
    \begin{steps}
        \item Postup u vývaru je stejný jako u polévky písmenkové.
        \item V hrníčku rozšleháme vidličkou vejce, přidáme mléko, koření a zašleháváme postupně polohrubou mouku, dokud nevznikne těstíčko. 
        \item Na kapání má být husté tak, aby těsto šlo přes struhadlo stěrkou protlačit. Protlačujeme rovnou do vroucí polévky a povaříme pár minut.
        \item Děláme-li nočky, těstíčko uděláme trochu hustší a z hrníčku vykrajujeme lžičkou noky.
    \end{steps}
\end{recipe}

\begin{recipe}[Fuchsia]{Česnečka}
    \begin{ingredients}
        \item \parbox[t]{\linewidth}{\raggedright \liquid \textbf{2 l}\\ Vývaru}
        \item \parbox[t]{\linewidth}{\raggedright \piece \textbf{3 plátky}\\ Tvrdšího chleba}
        \item \parbox[t]{\linewidth}{\raggedright \spoon \textbf{3 lžíce}\\ Sádla nebo oleje}
        \item \parbox[t]{\linewidth}{\raggedright \piece \textbf{5 stroužků}\\ Česneku}
        \item \parbox[t]{\linewidth}{\raggedright \weight \textbf{5 dkg}\\ Tvrdého sýra (eidam)}
        \item \parbox[t]{\linewidth}{\raggedright \weight \textbf{10 dkg}\\ Šunky}
        \item \parbox[t]{\linewidth}{\raggedright \indefinite \textbf{špetka}\\ Pepře a bylinek}
    \end{ingredients}
    \begin{steps}
        \item Přivedeme k varu vývar. Stáhneme z plotny a přidáme do něj 3 stroužky česneku, pepř a bylinky. 
        \item Na pánvi opečeme na sádle nebo oleji na kostičky nakrájený chléb. Po opečení vypneme oheň a do chleba vmícháme 2 stroužky česneku. 
        \item Do talíře dáme trochu nakrájeného sýra a šunky, zalijeme vývarem a na stůl dáme v misce opečený chléb.
    \end{steps}
\end{recipe}

\begin{recipe}[Plum]{Kouličková polévka od babči Libči}
    \begin{ingredients}
        \item \parbox[t]{\linewidth}{\raggedright \piece \textbf{1 ks}\\ Cibule}
        \item \parbox[t]{\linewidth}{\raggedright \spoon \textbf{1 PL}\\ Sádla}
        \item \parbox[t]{\linewidth}{\raggedright \piece \textbf{4 kostky}\\ Droždí}
        \item \parbox[t]{\linewidth}{\raggedright \piece \textbf{9 stroužků}\\ Česneku}
        \item \parbox[t]{\linewidth}{\raggedright \piece \textbf{1 ks}\\ Petržel a celer (menší)}
        \item \parbox[t]{\linewidth}{\raggedright \piece \textbf{2 ks}\\ Mrkve}
        \item \parbox[t]{\linewidth}{\raggedright \piece \textbf{3 hrnky}\\ Strouhanky}
        \item \parbox[t]{\linewidth}{\raggedright \piece \textbf{3 ks}\\ Vejce}
        \item \parbox[t]{\linewidth}{\raggedright \liquid \textbf{3 l}\\ Vývaru}
    \end{ingredients}
    \begin{steps}
        \item Na sádle necháme zesklovatět nadrobno nakrájenou cibuli. Přidáme droždí a na mírném plamenu necháme rozpustit. 
        \item Dáme vařit vývar s nadrobno nakrájenou zeleninou. 
        \item Do větší misky dáme strouhanku, osolíme, opepříme, přidáme prolisovaný česnek, majoránku, celá vejce a vychladlé droždí. 
        \item Přidáme trochu mléka nebo vývaru a vypracujeme tužší hmotu, ze které tvarujeme kouličky. 
        \item Kouličky vhazujeme do vroucího vývaru a vaříme, dokud nevyplavou nahoru.
    \end{steps}
\end{recipe}

\begin{recipe}[OliveGreen]{Bramboračka od babičky Jižňácké}
    \begin{ingredients}
        \item \parbox[t]{\linewidth}{\raggedright \liquid \textbf{2 l}\\ Vody}
        \item \parbox[t]{\linewidth}{\raggedright \piece \textbf{4 ks}\\ Menší brambory}
        \item \parbox[t]{\linewidth}{\raggedright \weight \textbf{300 g}\\ Zeleniny (mrkev, celer, květák...)}
        \item \parbox[t]{\linewidth}{\raggedright \piece \textbf{1 hrnek}\\ Sušených hub}
        \item \parbox[t]{\linewidth}{\raggedright \piece \textbf{2 kostky}\\ Masoxu}
        \item \parbox[t]{\linewidth}{\raggedright \indefinite \textbf{1 KL}\\ Majoránky a koření}
        \item \parbox[t]{\linewidth}{\raggedright \piece \textbf{2 stroužky}\\ Česneku}
        \item \parbox[t]{\linewidth}{\raggedright \spoon \textbf{dle potřeby}\\ Jíška}
    \end{ingredients}
    \begin{steps}
        \item Houby dám do hrnku, zaliji studenou vodou a nechám nabobtnat. 
        \item Dám do hrnce vodu, přidám na kostičky nakrájené brambory a zeleninu. Přidám masox nebo sůl a okořením. 
        \item Z hub sliji vodu a přidám je do hrnce. Přivedu k varu. 
        \item Přidám rozředěnou jíšku a vařím, dokud zelenina, brambory a houby nejsou měkké. Stáhnu z plotny a přidám drcený česnek.
    \end{steps}
\end{recipe}

\begin{recipe}[BurntOrange]{Cibulačka}
    \begin{ingredients}
        \item \parbox[t]{\linewidth}{\raggedright \piece \textbf{4 ks}\\ Velké cibule}
        \item \parbox[t]{\linewidth}{\raggedright \liquid \textbf{1,5 l}\\ Hovězího vývaru}
        \item \parbox[t]{\linewidth}{\raggedright \liquid \textbf{dle potřeby}\\ Olivový olej}
        \item \parbox[t]{\linewidth}{\raggedright \weight \textbf{100 g}\\ Parmazánu}
        \item \parbox[t]{\linewidth}{\raggedright \piece \textbf{2 ks}\\ Starší rohlíky}
        \item \parbox[t]{\linewidth}{\raggedright \piece \textbf{2 ks}\\ Vejce}
    \end{ingredients}
    \begin{steps}
        \item Cibule nakrájíme na půlkolečka a osmažíme dozlatova na oleji. 
        \item Zalijeme vývarem, přidáme libeček, sůl a pepř a 5-10 minut povaříme. Nakonec přidáme čerstvou petrželku. 
        \item Rohlíky nakrájíme na kolečka, namáčíme v rozšlehaném osoleném vejci a smažíme dozlatova. 
        \item Polévka se podává s těmito „pavézkami" a posypaná parmazánem.
    \end{steps}
\end{recipe}

\begin{recipe}[RawSienna]{Vývar}
    \begin{ingredients}
        \item \parbox[t]{\linewidth}{\raggedright \weight \textbf{300 g}\\ Masa s kostí}
        \item \parbox[t]{\linewidth}{\raggedright \liquid \textbf{2 l}\\ Vody}
        \item \parbox[t]{\linewidth}{\raggedright \piece \textbf{1 ks}\\ Cibule}
        \item \parbox[t]{\linewidth}{\raggedright \piece \textbf{2 ks}\\ Bobkové listy}
        \item \parbox[t]{\linewidth}{\raggedright \piece \textbf{7 kuliček}\\ Nového koření a pepře}
        \item \parbox[t]{\linewidth}{\raggedright \indefinite \textbf{dle chuti}\\ Sůl}
    \end{ingredients}
    \begin{steps}
        \item Do papiňáku dáme vodu, vložíme maso a kosti, přidáme koření. 
        \item Z cibule sloupneme jen ošklivou slupku, rozkrojíme napůl a přidáme k masu. Papiňák zavřeme a vaříme. 
        \item Délka vaření závisí na mase: Hovězí +/- 2 hodiny, vepřové +/- 1,5 hodiny, kuřecí +/- hodinu. 
        \item Po uvaření vývar přecedíme. Maso obereme a nakrájíme do polévky.
    \end{steps}
\end{recipe}

\begin{recipe}[BrickRed]{Směs na topinky}
    \begin{ingredients}
        \item \parbox[t]{\linewidth}{\raggedright \piece \textbf{1 ks}\\ Cibule}
        \item \parbox[t]{\linewidth}{\raggedright \weight \textbf{3 ks}\\ Kuřecí prsa}
        \item \parbox[t]{\linewidth}{\raggedright \liquid \textbf{1 ks}\\ Rajský protlak}
        \item \parbox[t]{\linewidth}{\raggedright \piece \textbf{5 kostek}\\ Cukru}
        \item \parbox[t]{\linewidth}{\raggedright \indefinite \textbf{1 balení}\\ Mražená zelenina}
        \item \parbox[t]{\linewidth}{\raggedright \piece \textbf{1 ks}\\ Červená paprika}
        \item \parbox[t]{\linewidth}{\raggedright \piece \textbf{2 stroužky}\\ Česneku}
        \item \parbox[t]{\linewidth}{\raggedright \spoon \textbf{1 PL}\\ Sójové omáčky}
        \item \parbox[t]{\linewidth}{\raggedright \piece \textbf{1 cm plátek}\\ Másla}
    \end{ingredients}
    \begin{steps}
        \item Cibuli oloupeme a nakrájíme nadrobno. Rozpustíme máslo, přidáme olej a osmažíme cibuli do zlatova. 
        \item Přidáme na drobné kostičky nakrájené kuřecí maso, osolíme a opepříme. Smažíme, dokud není maso skoro hotové. 
        \item Přidáme zeleninu, papriku, rajský protlak a cukr a dusíme do změknutí. 
        \item Dochutíme sójovou omáčkou, oreganem a drceným česnekem. Je-li omáčka řídká, zahustíme solamylem rozmíchaným ve vodě.
    \end{steps}
\end{recipe}

% ==========================================
% RECEPTY – UTOPENCI, OVAR A SLAVNOSTNÍ JÍDLA
% ==========================================

\begin{recipe}[IndianRed]{Utopenci od Bobana (Tomáš Horák z chaty v Řevnicích)}
    \begin{ingredients}
        \item \parbox[t]{\linewidth}{\raggedright \piece \textbf{15 ks}\\ Špekáčků}
        \item \parbox[t]{\linewidth}{\raggedright \liquid \textbf{0,5 l}\\ Vody a 0,2 l octa (nálev)}
        \item \parbox[t]{\linewidth}{\raggedright \piece \textbf{2 ks}\\ Cibule (na půlkolečka)}
        \item \parbox[t]{\linewidth}{\raggedright \piece \textbf{dle chuti}\\ Sterilované okurky a česnek}
        \item \parbox[t]{\linewidth}{\raggedright \indefinite \textbf{koření}\\ Sladká paprika, mletý pepř, feferonková drť}
        \item \parbox[t]{\linewidth}{\raggedright \piece \textbf{10 kuliček}\\ Celého pepře}
        \item \parbox[t]{\linewidth}{\raggedright \piece \textbf{3 kuličky}\\ Nového koření a 1–2 bobkové listy}
    \end{ingredients}
    \begin{steps}
        \item Svaříme 0,5 l vody a 0,2 l octa se solí, mletým pepřem, celým pepřem, novým kořením a bobkovým listem. Necháme vychladnout.
        \item Špekáčky oloupeme a rozkrojíme. Připravíme si cibuli na půlkolečka, sterilované okurky a koření.
        \item Do rozkrojeného špekáčku dáme špetku pepře, sladké papriky a feferonkové drtě (dle gusta), přidáme plátek okurky a cibuli.
        \item Naplněné špekáčky dáme natěsno do sklenice a prosypáváme cibulí, plátky okurek a případně beraními rohy.
        \item Zalijeme vychladlým lákem (teplý by na povrchu srazil tuk ze špekáčků) a necháme uležet minimálně 7 dní.
    \end{steps}
\end{recipe}

\begin{recipe}[Peru]{Ovar}
    \begin{ingredients}
        \item \parbox[t]{\linewidth}{\raggedright \weight \textbf{500 g}\\ Vepřového masa (nejraději plecko)}
        \item \parbox[t]{\linewidth}{\raggedright \weight \textbf{500 g}\\ Uzeného masa (např. krkovice s kostí)}
        \item \parbox[t]{\linewidth}{\raggedright \piece \textbf{10 kuliček}\\ Celého pepře}
        \item \parbox[t]{\linewidth}{\raggedright \piece \textbf{3 kuličky}\\ Nového koření}
        \item \parbox[t]{\linewidth}{\raggedright \piece \textbf{1 ks}\\ Bobkový list}
        \item \parbox[t]{\linewidth}{\raggedright \piece \textbf{1 větší}\\ Neoloupaná cibule}
    \end{ingredients}
    \begin{steps}
        \item Maso s kostí vložíme do hrnce a přidáme koření a rozkrojenou cibuli.
        \item Vaříme stejně jako vývar na polévku (viz recept Vývar), dokud není maso měkké.
        \item Podáváme s hořčicí, křenem a čerstvým chlebem.
    \end{steps}
\end{recipe}

\begin{recipe}[DarkGoldenrod]{Langoše}
    \begin{ingredients}
        \item \parbox[t]{\linewidth}{\raggedright \weight \textbf{500 g}\\ Hladké mouky}
        \item \parbox[t]{\linewidth}{\raggedright \piece \textbf{1 kostka}\\ Droždí}
        \item \parbox[t]{\linewidth}{\raggedright \liquid \textbf{4 dcl}\\ Mléka}
        \item \parbox[t]{\linewidth}{\raggedright \weight \textbf{1/2 KL}\\ Soli}
        \item \parbox[t]{\linewidth}{\raggedright \weight \textbf{1/2 KL}\\ Cukru}
    \end{ingredients}
    \begin{steps}
        \item Vlažné mléko smícháme s cukrem a rozdrobíme do něj droždí. Kvásek necháme asi 5 minut vzejít v teple.
        \item Mouku osolíme, prosejeme do mísy, přidáme kvásek a vypracujeme těsto. Necháme aspoň 30 minut kynout v teple.
        \item Těsto vyválíme na pomoučeném vále na 2 cm silný plát a sklenicí vykrojíme větší kolečka. Necháme ještě 10 minut vykynout.
        \item V pánvi rozpálíme vyšší vrstvu oleje. Nakynutá kolečka vytáhneme do stran a ihned smažíme dozlatova.
        \item Hotové langoše necháme okapat na savém papíru a doplníme vybranými surovinami (sýr, šunka, česnek, kečup, tatarka či marmeláda).
    \end{steps}
\end{recipe}

\begin{recipe}[SandyBrown]{Lívance}
    \begin{ingredients}
        \item \parbox[t]{\linewidth}{\raggedright \piece \textbf{1 kostka}\\ Droždí (na kvásek)}
        \item \parbox[t]{\linewidth}{\raggedright \spoon \textbf{2 PL}\\ Cukru (na kvásek)}
        \item \parbox[t]{\linewidth}{\raggedright \spoon \textbf{2 PL}\\ Hladké mouky (na kvásek)}
        \item \parbox[t]{\linewidth}{\raggedright \liquid \textbf{trochu}\\ Vlažného mléka (na kvásek)}
        \item \parbox[t]{\linewidth}{\raggedright \weight \textbf{500 g}\\ Hladké mouky}
        \item \parbox[t]{\linewidth}{\raggedright \piece \textbf{3 ks}\\ Vejce}
        \item \parbox[t]{\linewidth}{\raggedright \liquid \textbf{1 l}\\ Mléka}
        \item \parbox[t]{\linewidth}{\raggedright \indefinite \textbf{špetka}\\ Soli}
    \end{ingredients}
    \begin{steps}
        \item Mléko nalejeme do mísy, trošku osolíme, přidáme vejce a zašleháme.
        \item Přidáme hotový kvásek a hladkou mouku, pořádně rozmícháme kvedlačkou a necháme na teplém místě kynout.
        \item Po vykynutí naléváme těsto na vymaštěný lívanečník a opékáme z obou stran.
        \item Podáváme posypané moučkovým cukrem se skořicí, nebo potíráme marmeládou a zdobíme našlehaným tvarohem či pacholíkem.
    \end{steps}
\end{recipe}

\begin{recipe}[Tan]{Palačinky od Majky}
    \begin{ingredients}
        \item \parbox[t]{\linewidth}{\raggedright \liquid \textbf{1/2 l}\\ Mléka}
        \item \parbox[t]{\linewidth}{\raggedright \weight \textbf{20 dkg}\\ Hladké mouky}
        \item \parbox[t]{\linewidth}{\raggedright \piece \textbf{2 ks}\\ Vejce}
        \item \parbox[t]{\linewidth}{\raggedright \spoon \textbf{2 lžíce}\\ Olivového oleje}
    \end{ingredients}
    \begin{steps}
        \item Všechny suroviny smícháme na hladké těsto.
        \item Necháme odpočinout alespoň dvě hodiny, nejlépe přes noc.
        \item První palačinku pečeme na teflonu s pár kapkami oleje, další už bez oleje.
    \end{steps}
\end{recipe}

\begin{recipe}[Wheat]{Chléb obalený ve vajíčku}
    \begin{ingredients}
        \item \parbox[t]{\linewidth}{\raggedright \piece \textbf{1/2 ks}\\ Staršího chleba}
        \item \parbox[t]{\linewidth}{\raggedright \piece \textbf{4 ks}\\ Vejce}
        \item \parbox[t]{\linewidth}{\raggedright \indefinite \textbf{dle chuti}\\ Sůl, pepř, sušené bylinky}
        \item \parbox[t]{\linewidth}{\raggedright \liquid \textbf{dle potřeby}\\ Olej na smažení}
        \item \parbox[t]{\linewidth}{\raggedright \indefinite \textbf{na obložení}\\ Kečup, hořčice, cibule, kyselá okurka}
    \end{ingredients}
    \begin{steps}
        \item Starší chléb nakrájíme na krajíce, které přepůlíme.
        \item V misce rozšleháme vejce, osolíme, opepříme a přidáme rozemnuté sušené bylinky.
        \item Chléb ve vejcích obalíme a ihned vkládáme na pánev s rozehřátým olejem.
        \item Osmažíme a na ubrousku necháme odkapat přebytečný tuk. Ozdobíme dle chuti.
    \end{steps}
\end{recipe}

\begin{recipe}[Moccasin]{Smaženka}
    \begin{ingredients}
        \item \parbox[t]{\linewidth}{\raggedright \piece \textbf{5 krajíců}\\ Chleba}
        \item \parbox[t]{\linewidth}{\raggedright \piece \textbf{5 ks}\\ Vajec}
        \item \parbox[t]{\linewidth}{\raggedright \spoon \textbf{5 PL}\\ Strouhanky}
        \item \parbox[t]{\linewidth}{\raggedright \spoon \textbf{3 PL}\\ Mléka}
        \item \parbox[t]{\linewidth}{\raggedright \indefinite \textbf{dle chuti}\\ Sůl, pepř}
        \item \parbox[t]{\linewidth}{\raggedright \liquid \textbf{dle potřeby}\\ Olej na smažení}
        \item \parbox[t]{\linewidth}{\raggedright \indefinite \textbf{na oblohu}\\ Kečup, hořčice, cibule, kyselá okurka}
    \end{ingredients}
    \begin{steps}
        \item V misce rozšleháme vejce, přidáme mléko, strouhanku, osolíme a opepříme.
        \item Pánev potřeme olejem, rozehřejeme a sběračkou nalijeme část vaječné směsi tak, aby vznikla co nejtenčí smaženka.
        \item Po opečení z obou stran přendáme smaženku na krajíc chleba tak, aby polovina byla na chlebu a polovina volná.
        \item Polovinu na chlebu pomažeme hořčicí nebo kečupem, posypeme cibulí a okurkou a přiklopíme volnou polovinou.
    \end{steps}
\end{recipe}

\begin{recipe}[Beige]{Škrabánek}
    \begin{ingredients}
        \item \parbox[t]{\linewidth}{\raggedright \piece \textbf{zbytky}\\ Po obalování (vajíčko, mouka, strouhanka)}
        \item \parbox[t]{\linewidth}{\raggedright \indefinite \textbf{dle chuti}\\ Koření (paprika, bylinky, grilovací směs)}
        \item \parbox[t]{\linewidth}{\raggedright \liquid \textbf{trochu}\\ Mléka}
    \end{ingredients}
    \begin{steps}
        \item Při obalování v trojobalu zbývá v misce trocha vajíčka. Přidáme trošku koření, mléka a zahustíme hladkou moukou a strouhankou.
        \item Těstíčko musí být husté tak, aby se z něj dal vytvarovat škrabánek, ale ne moc husté, pak by byl příliš tuhý.
        \item Osmažíme ho z obou stran na pánvi.
    \end{steps}
\end{recipe}

\begin{recipe}[SkyBlue]{Rybí pomazánka od babičky Jižňácké}
    \begin{ingredients}
        \item \parbox[t]{\linewidth}{\raggedright \piece \textbf{250 g}\\ Másla (1 kostka)}
        \item \parbox[t]{\linewidth}{\raggedright \piece \textbf{3 plechovky}\\ Rybiček (sardinky od Oskara ve žlutém obalu)}
        \item \parbox[t]{\linewidth}{\raggedright \piece \textbf{1 plechovka}\\ Rybiček s feferonkou (oranžový obal)}
        \item \parbox[t]{\linewidth}{\raggedright \piece \textbf{4 ks}\\ Větší sladkokyselé okurky}
        \item \parbox[t]{\linewidth}{\raggedright \piece \textbf{1 ks}\\ Cibule}
        \item \parbox[t]{\linewidth}{\raggedright \indefinite \textbf{špetka}\\ Sladké papriky, sůl a pepř}
    \end{ingredients}
    \begin{steps}
        \item Máslo ve vyšší misce našleháme elektrickým šlehačem.
        \item Přidáme všechny rybičky a část oleje z plechovky a šleháme dále.
        \item Osolíme, opepříme a přidáme mletou papriku.
        \item Na závěr vmícháme nadrobno nakrájenou cibuli a okurku. Mažeme na pečivo a zdobíme cibulkou nebo okurkou.
    \end{steps}
\end{recipe}

\begin{recipe}[LightSeaGreen]{Celerová pomazánka}
    \begin{ingredients}
        \item \parbox[t]{\linewidth}{\raggedright \piece \textbf{1/4 ks}\\ Středního celeru}
        \item \parbox[t]{\linewidth}{\raggedright \piece \textbf{1 ks}\\ Jablko}
        \item \parbox[t]{\linewidth}{\raggedright \piece \textbf{3 ks}\\ Malé kvašáky}
        \item \parbox[t]{\linewidth}{\raggedright \piece \textbf{1/2 ks}\\ Červené cibule}
        \item \parbox[t]{\linewidth}{\raggedright \indefinite \textbf{hrst}\\ Nakrájené pažitky}
        \item \parbox[t]{\linewidth}{\raggedright \spoon \textbf{3 PL}\\ Majolky}
        \item \parbox[t]{\linewidth}{\raggedright \liquid \textbf{trocha}\\ Citronové šťávy}
        \item \parbox[t]{\linewidth}{\raggedright \indefinite \textbf{dle chuti}\\ Pepř}
    \end{ingredients}
    \begin{steps}
        \item Celer a jablko oloupeme a nastrouháme na drobných slzičkách.
        \item Přidáme najemno nakrájenou cibulku, okurku a pažitku. Zakápneme citrónem a lehce opepříme.
        \item Pořádně promícháme, přidáme majolku a znovu promícháme.
        \item Podáváme namazané na rohlíku nebo vece.
    \end{steps}
\end{recipe}

% ==========================================
% RECEPTY – SLADKÉ PEČENÍ A DEZERTY
% ==========================================

\begin{recipe}[Chocolate]{Tvaroh našlehaný s čokoládou od maminky}
    \begin{ingredients}
        \item \parbox[t]{\linewidth}{\raggedright \piece \textbf{1 vanička}\\ Tvarohu (polotučný/tučný)}
        \item \parbox[t]{\linewidth}{\raggedright \liquid \textbf{dle potřeby}\\ Mléko}
        \item \parbox[t]{\linewidth}{\raggedright \spoon \textbf{3 PL}\\ Moučkového cukru}
        \item \parbox[t]{\linewidth}{\raggedright \spoon \textbf{1/2 PL}\\ Kakaa}
    \end{ingredients}
    \begin{steps}
        \item Tvaroh dáme do vyšší misky a elektrickým šlehačem rozšleháme.
        \item Přidáme tolik mléka, aby vznikla pěna řidší konzistence, než je krupicová kaše. Rozdělíme do misek.
        \item V hrníčku smícháme cukr a kakao s trochou mléka, až vznikne tekutá čokoládová poleva.
        \item Polevu nalijeme na tvaroh. Pokud je konzistence správná, čokoláda klesne na dno.
    \end{steps}
\end{recipe}

\begin{recipe}[Goldenrod]{Listové tyčinky}
    \begin{ingredients}
        \item \parbox[t]{\linewidth}{\raggedright \piece \textbf{1 balíček}\\ Listového těsta}
        \item \parbox[t]{\linewidth}{\raggedright \piece \textbf{1 ks}\\ Vejce (na potření)}
        \item \parbox[t]{\linewidth}{\raggedright \indefinite \textbf{varianta 1}\\ Klobása a sýr}
        \item \parbox[t]{\linewidth}{\raggedright \indefinite \textbf{varianta 2}\\ Šunka a sýr}
        \item \parbox[t]{\linewidth}{\raggedright \indefinite \textbf{varianta 3}\\ Niva nebo tuňák}
    \end{ingredients}
    \begin{steps}
        \item Těsto rozválíme na velikost plechu. Položíme na pečicí papír.
        \item Pomažeme rozšlehaným vajíčkem a poklademe zvolenou ozdobou (sýr, šunka, koření).
        \item Rádýlkem rozkrájíme na tyčinky nebo čtverečky.
        \item Pečeme v rozpálené troubě na 200 °C asi 10 minut, pak snížíme na 170 °C a dopečeme dozlatova.
    \end{steps}
\end{recipe}

\begin{recipe}[RoyalPurple]{Prošívaná deka od Drissinky}
    \begin{ingredients}
        \item \parbox[t]{\linewidth}{\raggedright \piece \textbf{2 hrnky}\\ Polohrubé mouky}
        \item \parbox[t]{\linewidth}{\raggedright \piece \textbf{3 ks}\\ Vejce}
        \item \parbox[t]{\linewidth}{\raggedright \piece \textbf{1 ks}\\ Prášek do pečiva}
        \item \parbox[t]{\linewidth}{\raggedright \liquid \textbf{1/2 hrnku}\\ Mléka a 5 PL oleje}
        \item \parbox[t]{\linewidth}{\raggedright \piece \textbf{1 hrnek}\\ Cukru a 3 PL kakaa}
        \item \parbox[t]{\linewidth}{\raggedright \piece \textbf{2 ks}\\ Tvarohy (náplň)}
        \item \parbox[t]{\linewidth}{\raggedright \piece \textbf{1/2 ks}\\ Vanilkového pudinku (náplň)}
    \end{ingredients}
    \begin{steps}
        \item Smícháme suroviny na náplň (tvaroh, žloutky, pudink, cukr) a dáme do zdobicího sáčku.
        \item Smícháme přísady na tmavé těsto a nalijeme na vymazaný a vysypaný plech.
        \item Tvarohovou náplní naneseme na těsto pruhy do mřížky.
        \item Pečeme v předehřáté troubě, dokud těsto nevyběhne kolem tvarohu a nevytvoří efekt „deky".
    \end{steps}
\end{recipe}

\begin{recipe}[Plum]{Koláč tety Boženy od babičky Prosecké}
    \begin{ingredients}
        \item \parbox[t]{\linewidth}{\raggedright \piece \textbf{2 hrnky}\\ Polohrubé mouky}
        \item \parbox[t]{\linewidth}{\raggedright \piece \textbf{1 hrnek}\\ Cukru a 1 ks vejce}
        \item \parbox[t]{\linewidth}{\raggedright \liquid \textbf{1/2 hrnku}\\ Oleje a 1 hrnek mléka}
        \item \parbox[t]{\linewidth}{\raggedright \piece \textbf{1/2 ks}\\ Prášku do pečiva}
        \item \parbox[t]{\linewidth}{\raggedright \indefinite \textbf{dle sezóny}\\ Ovoce (švestky, meruňky)}
        \item \parbox[t]{\linewidth}{\raggedright \piece \textbf{drobenka}\\ Máslo, mouka, cukr}
    \end{ingredients}
    \begin{steps}
        \item Nejdříve si připravíme drobenku rozemnutím másla s cukrem a moukou.
        \item Všechny přísady na těsto smícháme a rozetřeme na vymazaný a vysypaný plech.
        \item Poklademe ovocem a hustě zasypeme drobenkou.
        \item Pečeme v předehřáté troubě dozlatova.
    \end{steps}
\end{recipe}

\begin{recipe}[Sienna]{Čoko muffiny}
    \begin{ingredients}
        \item \parbox[t]{\linewidth}{\raggedright \weight \textbf{330 g}\\ Polohrubé mouky}
        \item \parbox[t]{\linewidth}{\raggedright \piece \textbf{1 ks}\\ Prášek do pečiva}
        \item \parbox[t]{\linewidth}{\raggedright \spoon \textbf{3 PL}\\ Kakaa a 150 g cukru}
        \item \parbox[t]{\linewidth}{\raggedright \piece \textbf{3 ks}\\ Vejce}
        \item \parbox[t]{\linewidth}{\raggedright \liquid \textbf{270 ml}\\ Mléka a 188 g másla}
        \item \parbox[t]{\linewidth}{\raggedright \weight \textbf{150 g}\\ Nasekané čokolády}
    \end{ingredients}
    \begin{steps}
        \item Máslo rozehřejeme a v misce smícháme se všemi ostatními surovinami.
        \item Nakonec vmícháme nasekanou čokoládu a ořechy.
        \item Plníme do formiček na muffiny (silikonové nemusíme mazat).
        \item Pečeme na 180 °C přibližně 20–25 minut.
    \end{steps}
\end{recipe}

\begin{recipe}[DarkOrchid]{Trojbarevný koláč od Blanky Černé z VZLÚ}
    \begin{ingredients}
        \item \parbox[t]{\linewidth}{\raggedright \piece \textbf{1. těsto}\\ 2 hrnky mouky, 1/2 másla, kakao}
        \item \parbox[t]{\linewidth}{\raggedright \piece \textbf{2. těsto}\\ 3 tvarohy, cukr, 1 vejce}
        \item \parbox[t]{\linewidth}{\raggedright \piece \textbf{3. těsto}\\ 3 vejce, 1 hrnek mouky, 1/2 hrnku oleje}
    \end{ingredients}
    \begin{steps}
        \item První těsto zpracujeme prsty jako drobenku a napěchujeme na dno plechu.
        \item Na něj opatrně rozetřeme tvarohovou náplň (druhé těsto).
        \item Třetí těsto vyšleháme (vejce s cukrem do pěny, pak zbytek) a nalijeme navrch.
        \item Pečeme v troubě, dokud horní piškotová vrstva není hotová.
    \end{steps}
\end{recipe}

\begin{recipe}[BurntOrange]{Bábovka ze šlehačky od tety Aleny}
    \begin{ingredients}
        \item \parbox[t]{\linewidth}{\raggedright \piece \textbf{1/2 l}\\ Polohrubé mouky}
        \item \parbox[t]{\linewidth}{\raggedright \piece \textbf{1/4 l}\\ Cukru krystal}
        \item \parbox[t]{\linewidth}{\raggedright \piece \textbf{3 ks}\\ Vejce}
        \item \parbox[t]{\linewidth}{\raggedright \liquid \textbf{1/4 l}\\ Šlehačky (33\%)}
        \item \parbox[t]{\linewidth}{\raggedright \piece \textbf{1 ks}\\ Prášek do pečiva}
        \item \parbox[t]{\linewidth}{\raggedright \spoon \textbf{3 PL}\\ Kakaa (do půlky těsta)}
    \end{ingredients}
    \begin{steps}
        \item Všechny suroviny (kromě kakaa) smícháme v řidší těsto.
        \item Formu vymažeme a vysypeme hrubou moukou.
        \item Do formy nalijeme 2/3 těsta, zbytek obarvíme kakaem a dolijeme.
        \item Pečeme na 200 °C asi 10 min, pak stáhneme na 175 °C a dopékáme (celkem cca 45 min).
    \end{steps}
\end{recipe}

\begin{recipe}[DimGray]{Vláčný makovec}
    \begin{ingredients}
        \item \parbox[t]{\linewidth}{\raggedright \weight \textbf{300 g}\\ Mletého máku}
        \item \parbox[t]{\linewidth}{\raggedright \weight \textbf{130 g}\\ Polohrubé mouky}
        \item \parbox[t]{\linewidth}{\raggedright \liquid \textbf{200 ml}\\ Mléka a 200 g jogurtu}
        \item \parbox[t]{\linewidth}{\raggedright \piece \textbf{2 ks}\\ Vejce a 120 g cukru}
        \item \parbox[t]{\linewidth}{\raggedright \spoon \textbf{2 PL}\\ Džemu a 2 PL kakaa}
        \item \parbox[t]{\linewidth}{\raggedright \indefinite \textbf{citron}\\ Kůra a lžička sody}
    \end{ingredients}
    \begin{steps}
        \item Změklé máslo utřete s cukrem, solí, skořicí a kakaem.
        \item Zašlehejte vejce, džem, jogurt a citronovou kůru.
        \item Nakonec vmíchejte mák, mouku se sodou a mléko.
        \item Pečte v papírem vyložené formě na 170 °C asi 50 minut. Po vychladnutí můžete polít citronovou polevou.
    \end{steps}
\end{recipe}

\begin{recipe}[Goldenrod]{Bábovka s vaječným likérem}
    \begin{ingredients}
        \item \parbox[t]{\linewidth}{\raggedright \piece \textbf{5 ks}\\ Vajec a 200 g cukru}
        \item \parbox[t]{\linewidth}{\raggedright \liquid \textbf{250 ml}\\ Vaječného likéru}
        \item \parbox[t]{\linewidth}{\raggedright \liquid \textbf{250 ml}\\ Oleje}
        \item \parbox[t]{\linewidth}{\raggedright \weight \textbf{300 g}\\ Polohrubé mouky}
        \item \parbox[t]{\linewidth}{\raggedright \piece \textbf{1 ks}\\ Prášek do pečiva}
        \item \parbox[t]{\linewidth}{\raggedright \spoon \textbf{3 PL}\\ Kakaa (do části těsta)}
    \end{ingredients}
    \begin{steps}
        \item Vejce ušleháme s cukrem do pěny. Tenkým pramínkem zašleháme likér a olej.
        \item Vařečkou lehce vmícháme mouku s práškem do pečiva.
        \item Do vymazané a vysypané formy nalijeme polovinu těsta, druhou obarvíme kakaem a dolijeme.
        \item Pečeme na 200 °C (30 min), pak snížíme na 180 °C a dopečeme.
    \end{steps}
\end{recipe}

\begin{recipe}[BrickRed]{Perník s marmeládou a čokoládou od Martiny}
    \begin{ingredients}
        \item \parbox[t]{\linewidth}{\raggedright \piece \textbf{2 hrnky}\\ Polohrubé mouky}
        \item \parbox[t]{\linewidth}{\raggedright \liquid \textbf{1 hrnek}\\ Mléka a 1 hrnek oleje}
        \item \parbox[t]{\linewidth}{\raggedright \piece \textbf{1 hrnek}\\ Cukru a 2 vejce}
        \item \parbox[t]{\linewidth}{\raggedright \piece \textbf{1 ks}\\ Prášek do perníku}
        \item \parbox[t]{\linewidth}{\raggedright \spoon \textbf{2 PL}\\ Kakaa}
        \item \parbox[t]{\linewidth}{\raggedright \indefinite \textbf{poleva}\\ Čokoláda na vaření a tuk}
    \end{ingredients}
    \begin{steps}
        \item Vejce s cukrem šleháme do husté pěny (klidně 10 minut).
        \item Zašleháme olej. Mouku smícháme s práškem a kakaem a střídavě s mlékem vmícháme do směsi.
        \item Pečeme na 170 °C. Po vychladnutí rozřízneme, namažeme marmeládou a polijeme čokoládou rozpuštěnou ve vodní lázni.
    \end{steps}
\end{recipe}

\begin{recipe}[RoyalBlue]{Tachtlárna od babči Libči}
    \begin{ingredients}
        \item \parbox[t]{\linewidth}{\raggedright \piece \textbf{9 ks}\\ Vajec a 10 PL cukru}
        \item \parbox[t]{\linewidth}{\raggedright \weight \textbf{12 PL}\\ Polohrubé mouky}
        \item \parbox[t]{\linewidth}{\raggedright \liquid \textbf{1 dcl}\\ Oleje a 3/4 prdopeče}
        \item \parbox[t]{\linewidth}{\raggedright \piece \textbf{krém}\\ Máslo, tvaroh, cukr}
        \item \parbox[t]{\linewidth}{\raggedright \indefinite \textbf{ovoce}\\ Broskve, mandarinky}
        \item \parbox[t]{\linewidth}{\raggedright \liquid \textbf{poleva}\\ Šťáva z kompotu a pudink}
    \end{ingredients}
    \begin{steps}
        \item Upečeme piškot z vajec, cukru, oleje a mouky. Necháme vychladnout.
        \item Vyšleháme máslo s cukrem, vmícháme tvaroh a natřeme na korpus. Poklademe ovocem.
        \item Z kompotové šťávy (doplněné vodou na 800 ml) a pudinků uvaříme husté želé.
        \item Vlažnou polevu nalijeme na ovoce a necháme v chladu ztuhnout.
    \end{steps}
\end{recipe}

\begin{recipe}[DeepPink]{Řezy z tyčinky Margot od Simony Tobiškové starší}
    \begin{ingredients}
        \item \parbox[t]{\linewidth}{\raggedright \piece \textbf{2 hrnky}\\ Polohrubé mouky}
        \item \parbox[t]{\linewidth}{\raggedright \piece \textbf{2 ks}\\ Vejce a 1 hrnek cukru}
        \item \parbox[t]{\linewidth}{\raggedright \liquid \textbf{1 kelímek}\\ Šlehačky a 3 PL oleje}
        \item \parbox[t]{\linewidth}{\raggedright \piece \textbf{1 ks}\\ Velká Margotka (strouhaná)}
        \item \parbox[t]{\linewidth}{\raggedright \piece \textbf{1 ks}\\ Prášek do pečiva}
    \end{ingredients}
    \begin{steps}
        \item Nastrouháme Margotku a smícháme ji se všemi ostatními surovinami.
        \item Pokud je těsto moc husté, kápneme trochu mléka.
        \item Nalijeme na vymazaný plech a upečeme.
        \item Vychladlou buchtu polijeme čokoládovou polevou.
    \end{steps}
\end{recipe}

\begin{recipe}[Sienna]{Sokratovy perníčky od Honzy Šediny}
    \begin{ingredients}
        \item \parbox[t]{\linewidth}{\raggedright \weight \textbf{400 g}\\ Hladké mouky}
        \item \parbox[t]{\linewidth}{\raggedright \weight \textbf{150 g}\\ Moučkového cukru}
        \item \parbox[t]{\linewidth}{\raggedright \piece \textbf{2 ks}\\ Vejce}
        \item \parbox[t]{\linewidth}{\raggedright \spoon \textbf{2 PL}\\ Medu a 50 g másla}
        \item \parbox[t]{\linewidth}{\raggedright \spoon \textbf{1-2 KL}\\ Jedlé sody}
        \item \parbox[t]{\linewidth}{\raggedright \indefinite \textbf{koření}\\ Skořice, hřebíček, badyán}
    \end{ingredients}
    \begin{steps}
        \item Máslo s medem rozehřejeme. Vypracujeme těsto a necháme ho uležet.
        \item Vyválíme placku (0,5 cm) a vykrajujeme tvary.
        \item Pečeme na 180 °C asi 4–8 minut.
        \item Ihned po vytažení z trouby potřeme žloutkem pro krásný lesk.
    \end{steps}
\end{recipe}

% ==========================================
% RECEPTY – JABLKA, OŘECHY A ZÁVINY
% ==========================================

\begin{recipe}[Goldenrod]{Štrůdl}
    \begin{ingredients}
        \item \parbox[t]{\linewidth}{\raggedright \weight \textbf{500 g}\\ Hladké mouky}
        \item \parbox[t]{\linewidth}{\raggedright \weight \textbf{250 g}\\ Hery}
        \item \parbox[t]{\linewidth}{\raggedright \piece \textbf{1 kelímek}\\ Bílého jogurtu}
        \item \parbox[t]{\linewidth}{\raggedright \piece \textbf{1 ks}\\ Vejce (do těsta + na potření)}
        \item \parbox[t]{\linewidth}{\raggedright \indefinite \textbf{náplň}\\ Jablka, skořice, cukr}
        \item \parbox[t]{\linewidth}{\raggedright \indefinite \textbf{extra}\\ Ořechy, rozinky v rumu, strouhanka}
    \end{ingredients}
    \begin{steps}
        \item Z mouky, Hery, jogurtu a vejce vypracujeme těsto a necháme v lednici odležet.
        \item Rozdělíme na dvě části, vyválíme dotenka. Posypeme strouhankou (aby nepromáčela těsto).
        \item Poklademe nastrouhanými jablky, cukrem, skořicí, ořechy a rozinkami.
        \item Zavineme, okraje založíme dospod. Potřeme vejcem a propícháme vidličkou.
        \item Pečeme na 180 °C dozlatova.
    \end{steps}
\end{recipe}

\begin{recipe}[ForestGreen]{Netradiční tvarohový závin s ovocem}
    \begin{ingredients}
        \item \parbox[t]{\linewidth}{\raggedright \weight \textbf{300 g}\\ Hladké mouky a 220 g tuku}
        \item \parbox[t]{\linewidth}{\raggedright \piece \textbf{2 ks}\\ Žloutky (do těsta)}
        \item \parbox[t]{\linewidth}{\raggedright \piece \textbf{500 g}\\ Měkkého tvarohu (náplň)}
        \item \parbox[t]{\linewidth}{\raggedright \piece \textbf{2 ks}\\ Žloutky (do náplně)}
        \item \parbox[t]{\linewidth}{\raggedright \indefinite \textbf{350 g}\\ Jahod nebo jiného ovoce}
        \item \parbox[t]{\linewidth}{\raggedright \indefinite \textbf{posyp}\\ Mandlové lupínky}
    \end{ingredients}
    \begin{steps}
        \item Vypracujeme hladké těsto z mouky, tuku a žloutků. Rozdělíme na poloviny.
        \item Pláty potřeme směsí tvarohu, cukru a žloutků. Poklademe ovocem.
        \item Zavineme, potřeme rozšlehaným vejcem a posypeme mandlemi.
        \item Upečeme v předehřáté troubě dozlatova.
    \end{steps}
\end{recipe}

\begin{recipe}[Sienna]{Křehký jablečník}
    \begin{ingredients}
        \item \parbox[t]{\linewidth}{\raggedright \weight \textbf{500 g}\\ Hladké mouky}
        \item \parbox[t]{\linewidth}{\raggedright \weight \textbf{250 g}\\ Másla (studené)}
        \item \parbox[t]{\linewidth}{\raggedright \piece \textbf{2 ks}\\ Vejce a 2 PL zakysané smetany}
        \item \parbox[t]{\linewidth}{\raggedright \weight \textbf{130 g}\\ Moučkového cukru}
        \item \parbox[t]{\linewidth}{\raggedright \weight \textbf{1 kg}\\ Jablek (nastrouhaných)}
        \item \parbox[t]{\linewidth}{\raggedright \indefinite \textbf{koření}\\ Skořice, hřebíček, badyán}
    \end{ingredients}
    \begin{steps}
        \item Z mouky, cukru, másla, vajec a smetany rychle vypracujeme těsto. Necháme 20 min v chladu.
        \item Rozdělíme na dvě části. První vyválíme (ideálně mezi dvěma papíry) a dáme na plech.
        \item Poklademe vymačkanými jablky s kořením. Přikryjeme druhým plátem.
        \item Vidličkou spojíme okraje a propícháme povrch. Pečeme na 180 °C (25–35 min).
    \end{steps}
\end{recipe}

\begin{recipe}[Chocolate]{Ořechové řezy od Jarušky Semenské z VZLÚ}
    \begin{ingredients}
        \item \parbox[t]{\linewidth}{\raggedright \weight \textbf{500 g}\\ Hladké mouky (těsto)}
        \item \parbox[t]{\linewidth}{\raggedright \weight \textbf{200 g}\\ Másla a 2 PL medu (těsto)}
        \item \parbox[t]{\linewidth}{\raggedright \piece \textbf{300 g}\\ Ořechů (směs nasekaná)}
        \item \parbox[t]{\linewidth}{\raggedright \weight \textbf{120 g}\\ Másla a 3 KL medu (ořech. směs)}
        \item \parbox[t]{\linewidth}{\raggedright \liquid \textbf{3 dl}\\ Mléka (krém)}
        \item \parbox[t]{\linewidth}{\raggedright \weight \textbf{200 g}\\ Másla (krém)}
    \end{ingredients}
    \begin{steps}
        \item Těsto rozdělíme na dva díly. První placku upečeme samostatně (15 min).
        \item Na druhou placku dáme ořechovou směs (osmaženou dozlatova na pánvi) a upečeme.
        \item Z mléka, mouky a solamylu uvaříme kaši. Po vychladnutí vyšleháme s máslem, cukrem a citronem v krém.
        \item První placku pomažeme krémem a přiklopíme ořechovou plackou. Necháme 24 h uležet v lednici.
    \end{steps}
\end{recipe}

\begin{recipe}[Maroon]{Italské tiramisu}
    \begin{ingredients}
        \item \parbox[t]{\linewidth}{\raggedright \weight \textbf{500 g}\\ Mascarpone}
        \item \parbox[t]{\linewidth}{\raggedright \piece \textbf{5 ks}\\ Vajec}
        \item \parbox[t]{\linewidth}{\raggedright \spoon \textbf{5 PL}\\ Třtinového cukru}
        \item \parbox[t]{\linewidth}{\raggedright \liquid \textbf{1 šálek}\\ Silné kávy}
        \item \parbox[t]{\linewidth}{\raggedright \liquid \textbf{4 PL}\\ Amaretta}
        \item \parbox[t]{\linewidth}{\raggedright \piece \textbf{2 bal.}\\ Cukrářských piškotů}
    \end{ingredients}
    \begin{steps}
        \item Žloutky ušleháme s cukrem, bílky zvlášť na sníh. Do žloutků zašleháme Mascarpone a pak sníh.
        \item Piškoty jen krátce namáčíme v kávě s likérem a skládáme do formy.
        \item Střídáme vrstvy piškotů a krému. Končíme krémem.
        \item Bohatě zasypeme kakaem a necháme aspoň 2 hodiny v chladu.
    \end{steps}
\end{recipe}

% ==========================================
% RECEPTY – KYNUTÁ TĚSTA A PLACKY
% ==========================================

\begin{recipe}[ForestGreen]{Kynuté – nekynuté těsto}
    \begin{ingredients}
        \item \parbox[t]{\linewidth}{\raggedright \weight \textbf{60 dkg}\\ Hladké mouky}
        \item \parbox[t]{\linewidth}{\raggedright \piece \textbf{1 kostka}\\ Droždí}
        \item \parbox[t]{\linewidth}{\raggedright \weight \textbf{1/2 ks}\\ Hery (nebo 2 dcl oleje)}
        \item \parbox[t]{\linewidth}{\raggedright \piece \textbf{2 ks}\\ Vejce}
        \item \parbox[t]{\linewidth}{\raggedright \piece \textbf{1 ks}\\ Prášek do pečiva}
        \item \parbox[t]{\linewidth}{\raggedright \liquid \textbf{1 dcl}\\ Mléka}
    \end{ingredients}
    \begin{steps}
        \item Na vále zpracujeme rukama všechny suroviny v hladké těsto.
        \item **Pozor!** Toto těsto se nenechává kynout. Musí se zpracovat ihned!
        \item Hodí se skvěle na pizzu, koláče nebo domácí buchty.
    \end{steps}
\end{recipe}

\begin{recipe}[RawSienna]{Škvarkové placky od babi Tobi}
    \begin{ingredients}
        \item \parbox[t]{\linewidth}{\raggedright \weight \textbf{1 kg}\\ Mouky (1/2 hl., 1/2 pol.)}
        \item \parbox[t]{\linewidth}{\raggedright \weight \textbf{40 dkg}\\ Sekaných škvarků}
        \item \parbox[t]{\linewidth}{\raggedright \piece \textbf{1 kostka}\\ Droždí a 3 vejce}
        \item \parbox[t]{\linewidth}{\raggedright \weight \textbf{10 dkg}\\ Sádla}
        \item \parbox[t]{\linewidth}{\raggedright \liquid \textbf{600 ml}\\ Mléka}
        \item \parbox[t]{\linewidth}{\raggedright \indefinite \textbf{koření}\\ Kmín, pepř, sůl}
    \end{ingredients}
    \begin{steps}
        \item Z trochy mléka, cukru a droždí uděláme kvásek.
        \item Smícháme mouku, vejce, sádlo, koření, kvásek a škvarky. Vypracujeme těsto a necháme vykynout.
        \item Tvoříme placky, na plechu je necháme ještě chvíli dojít, potřeme olejem a uděláme mřížku.
        \item Pečeme na 200--220 °C, dopékáme horkovzduchem.
    \end{steps}
\end{recipe}

\begin{recipe}[Orange]{Karásky aneb kanafásky od babi Tobi}
    \begin{ingredients}
        \item \parbox[t]{\linewidth}{\raggedright \weight \textbf{75 dkg}\\ Mouky (půl hladké, půl polohrubé)}
        \item \parbox[t]{\linewidth}{\raggedright \piece \textbf{1 ks}\\ Žloutek a 1 celé vejce}
        \item \parbox[t]{\linewidth}{\raggedright \weight \textbf{5 dkg}\\ Sádla}
        \item \parbox[t]{\linewidth}{\raggedright \liquid \textbf{1 dcl}\\ Oleje}
        \item \parbox[t]{\linewidth}{\raggedright \piece \textbf{1 kostka}\\ Droždí}
        \item \parbox[t]{\linewidth}{\raggedright \liquid \textbf{1/2 l}\\ Mléka}
        \item \parbox[t]{\linewidth}{\raggedright \spoon \textbf{1 PL}\\ Cukru (na kvásek)}
        \item \parbox[t]{\linewidth}{\raggedright \indefinite \textbf{dle chuti}\\ Sůl (trochu víc)}
    \end{ingredients}
    \begin{steps}
        \item Z cca 1/2 hrnku mléka uděláme vlažný kvásek s droždím, cukrem a lžící mouky.
        \item Smícháme všechny suroviny s kváskem a zaděláme těsto. Necháme vykynout.
        \item Upleteme ze tří pramenů kanafásky, potřeme bílkem a posypeme mákem, ořechy, semínky či solí.
        \item Pečeme na pečícím papíru na 200--220 °C, dopékáme horkovzduchem.
    \end{steps}
\end{recipe}

\begin{recipe}[NavyBlue]{Domácí chléb}
    \begin{ingredients}
        \item \parbox[t]{\linewidth}{\raggedright \weight \textbf{1 kg}\\ Hladké mouky chlebové}
        \item \parbox[t]{\linewidth}{\raggedright \liquid \textbf{6 dl}\\ Vlažné vody}
        \item \parbox[t]{\linewidth}{\raggedright \piece \textbf{1 kostka}\\ Droždí (42 g)}
        \item \parbox[t]{\linewidth}{\raggedright \liquid \textbf{2 lžíce}\\ Oleje}
        \item \parbox[t]{\linewidth}{\raggedright \spoon \textbf{2 lžičky}\\ Soli a 2 lžíce kmínu}
    \end{ingredients}
    \begin{steps}
        \item Uděláme kvásek z vody, droždí a cukru. Smícháme s moukou a ostatními surovinami.
        \item Vypracujeme těsto a necháme 45 min kynout.
        \item Vytvarujeme bochník na vymazaném plechu a necháme dalších 30 min kynout.
        \item Třikrát nařízneme, potřeme vejcem a pečeme 15 min na 225 °C, poté 25 min na 200 °C.
    \end{steps}
\end{recipe}

\begin{recipe}[DarkKhaki]{Bramborové placky}
    \begin{ingredients}
        \item \parbox[t]{\linewidth}{\raggedright \weight \textbf{500 g}\\ Brambor}
        \item \parbox[t]{\linewidth}{\raggedright \weight \textbf{100 g}\\ Polohrubé mouky}
        \item \parbox[t]{\linewidth}{\raggedright \piece \textbf{1 ks}\\ Celé vejce}
        \item \parbox[t]{\linewidth}{\raggedright \indefinite \textbf{dle potřeby}\\ Sádlo a sůl}
        \item \parbox[t]{\linewidth}{\raggedright \indefinite \textbf{hrst}\\ Zelené petrželky}
        \item \parbox[t]{\linewidth}{\raggedright \weight \textbf{20 dkg}\\ Šunky (nemusí být)}
    \end{ingredients}
    \begin{steps}
        \item Brambory uvaříme ve slupce cca 20 minut, oloupeme, necháme vychladnout a nastrouháme najemno.
        \item Přidáme sůl, vejce, mouku, trochu rozpuštěného sádla, nasekanou petrželku (případně šunku) a vypracujeme těsto.
        \item Vyválíme silnější placku, z které vykrájíme menší kolečka.
        \item Smažíme na oleji nebo rozpuštěném sádle. Podle chuti ještě osolíme bylinkovou solí.
    \end{steps}
\end{recipe}

% ==========================================
% RECEPTY – SALÁTY A ZELENINA
% ==========================================

\begin{recipe}[SpringGreen]{Salát okurkový}
    \begin{ingredients}
        \item \parbox[t]{\linewidth}{\raggedright \piece \textbf{1 ks}\\ Salátové okurky}
        \item \parbox[t]{\linewidth}{\raggedright \spoon \textbf{5 PL}\\ Cukru krystal}
        \item \parbox[t]{\linewidth}{\raggedright \spoon \textbf{1--1,5 PL}\\ Octa}
        \item \parbox[t]{\linewidth}{\raggedright \indefinite \textbf{dle chuti}\\ Pepř}
    \end{ingredients}
    \begin{steps}
        \item Okurku oloupeme a nastrouháme na struhadle na větších slzičkách.
        \item Zasypeme cukrem a necháme chvíli, aby pustila šťávu.
        \item Trochu opepříme a dochutíme octem. Podáváme vychlazené.
    \end{steps}
\end{recipe}

\begin{recipe}[DarkOrange]{Salát mrkvový}
    \begin{ingredients}
        \item \parbox[t]{\linewidth}{\raggedright \weight \textbf{500 g}\\ Mrkve}
        \item \parbox[t]{\linewidth}{\raggedright \spoon \textbf{6 PL}\\ Cukru krystal}
        \item \parbox[t]{\linewidth}{\raggedright \piece \textbf{1 ks}\\ Citrón}
    \end{ingredients}
    \begin{steps}
        \item Mrkev oškrábeme a nastrouháme na struhadle s menšími slzičkami.
        \item Posypeme krystalovým cukrem a promícháme.
        \item Necháme chvilku uležet a pak přidáme vodu a šťávu z citrónu podle chuti. Podáváme vychlazené.
    \end{steps}
\end{recipe}

\begin{recipe}[DarkSeaGreen]{Salát Coleslaw}
    \begin{ingredients}
        \item \parbox[t]{\linewidth}{\raggedright \piece \textbf{1 ks}\\ Menší hlávka bílého zelí}
        \item \parbox[t]{\linewidth}{\raggedright \piece \textbf{1 ks}\\ Cibule}
        \item \parbox[t]{\linewidth}{\raggedright \piece \textbf{2 ks}\\ Mrkve}
        \item \parbox[t]{\linewidth}{\raggedright \piece \textbf{1/4 ks}\\ Celeru}
        \item \parbox[t]{\linewidth}{\raggedright \spoon \textbf{1 PL}\\ Soli}
        \item \parbox[t]{\linewidth}{\raggedright \liquid \textbf{100 ml}\\ Smetany ke šlehání (zálivka)}
        \item \parbox[t]{\linewidth}{\raggedright \weight \textbf{200 g}\\ Majonézy (zálivka)}
        \item \parbox[t]{\linewidth}{\raggedright \spoon \textbf{2 PL}\\ Cukru krupice a 1 KL octa (zálivka)}
    \end{ingredients}
    \begin{steps}
        \item Zelí nakrouháme na tenké nudličky, přidáme lžíci soli a důkladně 20 minut mačkáme, aby pustilo co nejvíce šťávy. Tekutinu slijeme.
        \item Mrkev a celer nastrouháme nahrubo, cibuli nakrájíme najemno a vše promícháme.
        \item Zálivku připravíme smícháním majolky, smetany, cukru, octa a pepře.
        \item Salát se zálivkou promícháme a dochutíme. Podáváme uležené a vychlazené.
    \end{steps}
\end{recipe}

\begin{recipe}[MediumPurple]{Zelí bílé a červené}
    \begin{ingredients}
        \item \parbox[t]{\linewidth}{\raggedright \piece \textbf{1 sklenice}\\ Zelí (bílé sterilované, kysané nebo červené)}
        \item \parbox[t]{\linewidth}{\raggedright \piece \textbf{1 ks}\\ Cibule}
        \item \parbox[t]{\linewidth}{\raggedright \indefinite \textbf{kousek}\\ Sádla nebo trocha oleje}
        \item \parbox[t]{\linewidth}{\raggedright \spoon \textbf{1 PL}\\ Hladké mouky}
        \item \parbox[t]{\linewidth}{\raggedright \liquid \textbf{dle chuti}\\ Ocet a cukr na dochucení}
    \end{ingredients}
    \begin{steps}
        \item Cibuli nakrájíme nadrobno a na tuku osmahneme dozlatova.
        \item Zelí překrájíme na menší kousky, přidáme k cibuli, promícháme a podlijeme trochou vody. Dusíme pod pokličkou.
        \item Když je zelí skoro měkké, lehce osolíme, dochutíme cukrem a zaprášíme hladkou moukou. Za stálého míchání ještě chvíli podusíme.
        \item Kysané zelí přislazujeme více a ocet nedáváme. Sterilované bílé nebo červené zelí nesladíme tolik a dochucujeme octem.
    \end{steps}
\end{recipe}

\begin{recipe}[LimeGreen]{Špenát}
    \begin{ingredients}
        \item \parbox[t]{\linewidth}{\raggedright \piece \textbf{1 ks}\\ Špenát (protlak)}
        \item \parbox[t]{\linewidth}{\raggedright \piece \textbf{1 ks}\\ Cibule (nemusí být)}
        \item \parbox[t]{\linewidth}{\raggedright \indefinite \textbf{kousek}\\ Sádla nebo trocha oleje}
        \item \parbox[t]{\linewidth}{\raggedright \piece \textbf{5 stroužků}\\ Česneku}
        \item \parbox[t]{\linewidth}{\raggedright \spoon \textbf{1 PL}\\ Hladké mouky}
        \item \parbox[t]{\linewidth}{\raggedright \piece \textbf{1 ks}\\ Vejce (nemusí být)}
    \end{ingredients}
    \begin{steps}
        \item Cibuli nakrájíme nadrobno a na tuku osmahneme dozlatova.
        \item Špenát vložíme do kastrolu, lehce podlijeme, přiklopíme a často mícháme, dokud se nerozmrazí.
        \item Dusíme cca 10 minut, pak lehce osolíme, opepříme a zašleháme vejce.
        \item Zaprášíme moukou, přidáme prolisovaný česnek a ještě chvíli za stálého míchání podusíme.
    \end{steps}
\end{recipe}

% ==========================================
% RECEPTY – KNEDLÍKY, PŘÍLOHY A GULÁŠE
% ==========================================

\begin{recipe}[Khaki]{Jemný knedlík v ubrousku aneb bábovkové knedlíky od maminky}
    \begin{ingredients}
        \item \parbox[t]{\linewidth}{\raggedright \piece \textbf{2 ks}\\ Rohlíky}
        \item \parbox[t]{\linewidth}{\raggedright \spoon \textbf{2 lžíce}\\ Másla}
        \item \parbox[t]{\linewidth}{\raggedright \piece \textbf{3 ks}\\ Žloutky}
        \item \parbox[t]{\linewidth}{\raggedright \liquid \textbf{0,33 l}\\ Mléka}
        \item \parbox[t]{\linewidth}{\raggedright \indefinite \textbf{dle chuti}\\ Sůl}
        \item \parbox[t]{\linewidth}{\raggedright \weight \textbf{300 g}\\ Polohrubé mouky}
        \item \parbox[t]{\linewidth}{\raggedright \piece \textbf{3 ks}\\ Ušlehané bílky}
        \item \parbox[t]{\linewidth}{\raggedright \spoon \textbf{5 lžic}\\ Nasekaných zelených bylinek}
        \item \parbox[t]{\linewidth}{\raggedright \piece \textbf{5 stroužků}\\ Česneku}
    \end{ingredients}
    \begin{steps}
        \item Rohlíky nakrájíme na drobné kostičky, osmažíme na másle a necháme vychladnout.
        \item Žloutky rozmícháme v mléce, přidáme bylinky, česnek, sůl a za stálého míchání přisypáváme mouku.
        \item Vypracujeme řídké těsto, do kterého přidáme opečené rohlíky a lehce vmícháme sníh z bílků.
        \item Těsto nalijeme do velkého plátěného ubrousku namočeného ve studené vodě a potřeného tukem. Zavineme, svážeme motouzem a zavěsíme na vařečku do vroucí osolené vody. Vaříme přikryté 15 minut, pak přestřihneme první motouz a vaříme 15 minut, obrátíme a dovaříme ještě 15 minut. Vyjmeme a nakrájíme.
        \item Varianta: těsto nalijeme do vymazané bábovkové formy a vaříme ve vodní lázni 45 minut. Nebo nalijeme do vymazaných a vysypaných hrnků do 2/3, přikryjeme mikrotenem a dáme do mikrovlnky na plný výkon cca 8 minut.
    \end{steps}
\end{recipe}

\begin{recipe}[Coral]{Nádivka (ke kuřeti)}
    \begin{ingredients}
        \item \parbox[t]{\linewidth}{\raggedright \piece \textbf{6 ks}\\ Starších rohlíků}
        \item \parbox[t]{\linewidth}{\raggedright \weight \textbf{200 g}\\ Anglické slaniny}
        \item \parbox[t]{\linewidth}{\raggedright \weight \textbf{100 g}\\ Másla}
        \item \parbox[t]{\linewidth}{\raggedright \piece \textbf{6 ks}\\ Vajec}
        \item \parbox[t]{\linewidth}{\raggedright \liquid \textbf{100 ml}\\ Mléka}
        \item \parbox[t]{\linewidth}{\raggedright \indefinite \textbf{2 hrsti}\\ Zelených natí (kopřiva, petržel, celer)}
        \item \parbox[t]{\linewidth}{\raggedright \indefinite \textbf{dle chuti}\\ Muškátový oříšek}
        \item \parbox[t]{\linewidth}{\raggedright \piece \textbf{6 stroužků}\\ Česneku}
        \item \parbox[t]{\linewidth}{\raggedright \indefinite \textbf{dle chuti}\\ Sůl a pepř}
    \end{ingredients}
    \begin{steps}
        \item Rohlíky nakrájíme na kostičky a zalijeme mlékem.
        \item Přidáme žloutky, rozpuštěné máslo, nakrájenou slaninu, bylinky, strouhaný muškátový oříšek, drcený česnek, sůl a pepř.
        \item Důkladně promícháme a opatrně vmícháme z bílků ušlehaný sníh.
        \item Hmotou naplníme kuře a zašijeme. Zbylou nádivku upravíme do bochánku a dáme vedle kuřete, při pečení poléváme vypečenou šťávou.
        \item Nádivku můžeme péci i samostatně ve vymazaném a strouhankou vysypaném pekáčku.
    \end{steps}
\end{recipe}

\begin{recipe}[Gold]{Brkaše aneb bramborová kaše od maminky}
    \begin{ingredients}
        \item \parbox[t]{\linewidth}{\raggedright \weight \textbf{1 kg}\\ Brambor (dle počtu strávníků)}
        \item \parbox[t]{\linewidth}{\raggedright \liquid \textbf{1 miska}\\ Mléka}
        \item \parbox[t]{\linewidth}{\raggedright \piece \textbf{cca 2 cm}\\ Plátek másla}
        \item \parbox[t]{\linewidth}{\raggedright \indefinite \textbf{dle chuti}\\ Sůl}
    \end{ingredients}
    \begin{steps}
        \item Brambory oloupeme, nakrájíme a uvaříme v osolené vodě doměkka.
        \item Brambory scedíme, přidáme máslo a šťouchadlem rozšťoucháme.
        \item Elektrickým šlehačem šleháme tak dlouho, dokud nemáme žádné hrudky.
        \item Teprve když je hmota hladká, přidáváme po částech horké mléko. Dle chuti přisolit.
    \end{steps}
\end{recipe}

\begin{recipe}[Brown]{Buřtguláš}
    \begin{ingredients}
        \item \parbox[t]{\linewidth}{\raggedright \weight \textbf{1 kg}\\ Brambor}
        \item \parbox[t]{\linewidth}{\raggedright \weight \textbf{400 g}\\ Točeného salámu (nebo 4 ks buřtů či 500 g mletého masa)}
        \item \parbox[t]{\linewidth}{\raggedright \piece \textbf{2 ks}\\ Větší cibule}
        \item \parbox[t]{\linewidth}{\raggedright \spoon \textbf{1 PL}\\ Hladké mouky}
        \item \parbox[t]{\linewidth}{\raggedright \spoon \textbf{2 PL}\\ Mleté sladké papriky}
        \item \parbox[t]{\linewidth}{\raggedright \liquid \textbf{dle potřeby}\\ Olej na osmažení cibule}
        \item \parbox[t]{\linewidth}{\raggedright \indefinite \textbf{dle chuti}\\ Mletý kmín, sůl, pepř, majoránka}
        \item \parbox[t]{\linewidth}{\raggedright \piece \textbf{2 stroužky}\\ Česneku}
    \end{ingredients}
    \begin{steps}
        \item Cibuli oloupeme a nakrájíme nadrobno. Na troše oleje nebo sádla ji osmažíme.
        \item Přidáme na kostičky nakrájený salám a ještě chvíli opékáme. Pokud použijeme mleté maso, opékáme déle, dokud není skoro měkké.
        \item Vsypeme mletou papriku, promícháme a podlijeme vodou.
        \item Přidáme na kostičky nakrájené brambory, koření a dusíme do poloměkka.
        \item V hrníčku naplněném do poloviny vodou rozkvedláme mouku a vlijeme do guláše. Na závěr přidáme drcený česnek.
    \end{steps}
\end{recipe}

\begin{recipe}[Mahogany]{Guláš}
    \begin{ingredients}
        \item \parbox[t]{\linewidth}{\raggedright \weight \textbf{1 kg}\\ Masa (plec, kýta nebo hovězí zadní)}
        \item \parbox[t]{\linewidth}{\raggedright \weight \textbf{0,5 -- 1 kg}\\ Cibule}
        \item \parbox[t]{\linewidth}{\raggedright \spoon \textbf{2 PL}\\ Mleté sladké papriky}
        \item \parbox[t]{\linewidth}{\raggedright \piece \textbf{3 stroužky}\\ Česneku}
        \item \parbox[t]{\linewidth}{\raggedright \indefinite \textbf{koření}\\ Sůl, pepř, kmín, majoránka, chilli}
        \item \parbox[t]{\linewidth}{\raggedright \spoon \textbf{1--2 PL}\\ Hladké mouky (na zahuštění)}
    \end{ingredients}
    \begin{steps}
        \item Nadrobno nakrájenou cibuli osmažíme na sádle, přidáme maso na kostky a opékáme, dokud nepustí šťávu.
        \item Zaprášíme paprikou, podlijeme vodou, přidáme koření a dusíme doměkka.
        \item Zahustíme moukou rozmíchanou ve vodě a ještě 10 minut povaříme.
        \item Nakonec přidáme drcený česnek a majoránku.
    \end{steps}
\end{recipe}

% ==========================================
% RECEPTY – MASITÉ POKRMY
% ==========================================

\begin{recipe}[DarkOliveGreen]{Bramborový knedlík plněný uzeným}
    \begin{ingredients}
        \item \parbox[t]{\linewidth}{\raggedright \piece \textbf{6 ks}\\ Středně velkých brambor}
        \item \parbox[t]{\linewidth}{\raggedright \weight \textbf{200 g}\\ Hrubé mouky (podle potřeby i víc)}
        \item \parbox[t]{\linewidth}{\raggedright \weight \textbf{50 g}\\ Krupice}
        \item \parbox[t]{\linewidth}{\raggedright \piece \textbf{2 ks}\\ Vejce}
        \item \parbox[t]{\linewidth}{\raggedright \indefinite \textbf{dle chuti}\\ Sůl}
        \item \parbox[t]{\linewidth}{\raggedright \weight \textbf{30 dkg}\\ Vařeného uzeného masa (klidně i víc)}
        \item \parbox[t]{\linewidth}{\raggedright \piece \textbf{2 ks}\\ Cibule}
        \item \parbox[t]{\linewidth}{\raggedright \spoon \textbf{2 PL}\\ Sádla}
    \end{ingredients}
    \begin{steps}
        \item Brambory uvaříme ve slupce, po vychladnutí oloupeme a nastrouháme najemno.
        \item Přidáme vejce, sůl, krupici, zamícháme a přidáme hrubou mouku. Vypracujeme těsto.
        \item Z těsta děláme placky, na které dáváme na drobné kostičky nakrájené uzené. Zabalíme do kuliček.
        \item Lépe se pracuje navlhčenýma rukama, těsto se tolik nelepí. Vkládáme do vroucí osolené vody a po vyplavání vaříme cca 3 minuty.
        \item Varianta: místo domácího těsta použijeme pytlík bramborových knedlíků v prášku s trošku méně vody a jedním vejcem navíc.
        \item Podáváme se zelím a s na tuku opečenou nadrobno nakrájenou cibulkou.
    \end{steps}
\end{recipe}

\begin{recipe}[FireBrick]{Moravský vrabec}
    \begin{ingredients}
        \item \parbox[t]{\linewidth}{\raggedright \weight \textbf{1 kg}\\ Vepřového masa (plecko nebo kýta, můžeme přidat trochu bůčku)}
        \item \parbox[t]{\linewidth}{\raggedright \piece \textbf{3 ks}\\ Velké cibule}
        \item \parbox[t]{\linewidth}{\raggedright \indefinite \textbf{dle potřeby}\\ Sádlo}
        \item \parbox[t]{\linewidth}{\raggedright \piece \textbf{5 stroužků}\\ Česneku}
        \item \parbox[t]{\linewidth}{\raggedright \indefinite \textbf{dle chuti}\\ Mletý kmín, sůl, pepř, špetka chilli}
    \end{ingredients}
    \begin{steps}
        \item Maso pokrájíme na kostky, cibuli nakrájíme nahrubo.
        \item Vymažeme pekáček sádlem, vložíme maso a promícháme s cibulí, česnekem, solí, pepřem, chilli a kmínem.
        \item Pokud použijeme jen libové maso, přidáme navrch dvě lžíce sádla.
        \item Pečeme v troubě na 200 °C tak dlouho, dokud maso nepustí šťávu a vypečená šťáva se nevysmahne a maso se trošku „nepřipeče". Pečlivě hlídáme, aby nebylo hořké.
        \item Teprve pak podlijeme a pečeme, dokud není maso měkké. Podáváme s knedlíkem a dušeným zelím či špenátem.
    \end{steps}
\end{recipe}

\begin{recipe}[Tomato]{Pečené kuřátko}
    \begin{ingredients}
        \item \parbox[t]{\linewidth}{\raggedright \piece \textbf{1 ks}\\ Celé kuře nebo kuřecí kousky}
        \item \parbox[t]{\linewidth}{\raggedright \indefinite \textbf{dle chuti}\\ Koření (grilovací, na pečené kuře, gyros...)}
        \item \parbox[t]{\linewidth}{\raggedright \piece \textbf{kousek}\\ Másla}
    \end{ingredients}
    \begin{steps}
        \item Kuře očistíme a vetřeme do něj koření. Pokud koření obsahuje sůl, již nesolíme.
        \item Kuře můžeme naplnit nádivkou, dáme na něj kousek másla a vložíme do pekáčku.
        \item Podlijeme trochou vody a pečeme na 200 °C dozlatova.
        \item Po upečení naporcujeme a podáváme nejčastěji s bramborem.
    \end{steps}
\end{recipe}

\begin{recipe}[RosyBrown]{Karbanátky}
    \begin{ingredients}
        \item \parbox[t]{\linewidth}{\raggedright \weight \textbf{1 kg}\\ Mletého masa (nejlépe míchané vepřové a hovězí)}
        \item \parbox[t]{\linewidth}{\raggedright \piece \textbf{1 ks}\\ Menší cibule}
        \item \parbox[t]{\linewidth}{\raggedright \piece \textbf{4 stroužky}\\ Česneku}
        \item \parbox[t]{\linewidth}{\raggedright \piece \textbf{2 ks}\\ Vejce}
        \item \parbox[t]{\linewidth}{\raggedright \indefinite \textbf{dle chuti}\\ Sůl, pepř, mletý kmín, majoránka}
        \item \parbox[t]{\linewidth}{\raggedright \spoon \textbf{1 PL}\\ Hladké mouky}
        \item \parbox[t]{\linewidth}{\raggedright \indefinite \textbf{hrst}\\ Strouhanky na zahuštění}
        \item \parbox[t]{\linewidth}{\raggedright \indefinite \textbf{dle potřeby}\\ Trojobal a olej na smažení}
    \end{ingredients}
    \begin{steps}
        \item Do mísy dáme mleté maso, vajíčka, nadrobno nakrájenou cibuli, nadrcený česnek, koření a mouku.
        \item Pořádně promícháme a postupně zahušťujeme strouhankou. Můžeme přidat i trošku mléka.
        \item Uděláme placičky, které buď obalíme v trojobalu, jen ve strouhance, nebo je osmažíme jen tak.
        \item Smažíme na pánvi na oleji dozlatova.
    \end{steps}
\end{recipe}

\begin{recipe}[OrangeRed]{Čevabčiči}
    \begin{ingredients}
        \item \parbox[t]{\linewidth}{\raggedright \weight \textbf{1 kg}\\ Mletého masa}
        \item \parbox[t]{\linewidth}{\raggedright \spoon \textbf{2 KL}\\ Mleté sladké papriky}
        \item \parbox[t]{\linewidth}{\raggedright \indefinite \textbf{špetka}\\ Pálivé papriky}
        \item \parbox[t]{\linewidth}{\raggedright \piece \textbf{2 stroužky}\\ Česneku}
        \item \parbox[t]{\linewidth}{\raggedright \indefinite \textbf{dle chuti}\\ Sůl a pepř}
    \end{ingredients}
    \begin{steps}
        \item V míse smícháme všechny uvedené suroviny.
        \item Tvoříme malé válečky, které smažíme na oleji na pánvi.
        \item Podáváme s bramborem, cibulí a plnotučnou hořčicí.
    \end{steps}
\end{recipe}

\begin{recipe}[SeaGreen]{Holandský řízek}
    \begin{ingredients}
        \item \parbox[t]{\linewidth}{\raggedright \weight \textbf{600 g}\\ Mletého masa (má být bůček)}
        \item \parbox[t]{\linewidth}{\raggedright \liquid \textbf{1 dl}\\ Mléka}
        \item \parbox[t]{\linewidth}{\raggedright \weight \textbf{200 g}\\ Tvrdého sýra (eidam)}
        \item \parbox[t]{\linewidth}{\raggedright \indefinite \textbf{dle chuti}\\ Sůl a pepř}
        \item \parbox[t]{\linewidth}{\raggedright \indefinite \textbf{dle potřeby}\\ Trojobal}
    \end{ingredients}
    \begin{steps}
        \item Mleté maso promícháme s mlékem, solí a pepřem.
        \item Přidáme nastrouhaný sýr.
        \item Tvoříme placičky, které obalujeme v trojobalu a smažíme na rozpáleném oleji na pánvi.
        \item Podáváme nejčastěji s bramborovou kaší.
    \end{steps}
\end{recipe}

\begin{recipe}[Olive]{Omáčka na špagety od maminky}
    \begin{ingredients}
        \item \parbox[t]{\linewidth}{\raggedright \piece \textbf{2 ks}\\ Cibule}
        \item \parbox[t]{\linewidth}{\raggedright \spoon \textbf{1 PL}\\ Sádla nebo 2 PL oleje}
        \item \parbox[t]{\linewidth}{\raggedright \weight \textbf{300 g}\\ Anglické slaniny nebo šunky}
        \item \parbox[t]{\linewidth}{\raggedright \piece \textbf{1 ks}\\ Rajský protlak}
        \item \parbox[t]{\linewidth}{\raggedright \piece \textbf{3 kostky}\\ Cukru (podle chuti)}
        \item \parbox[t]{\linewidth}{\raggedright \indefinite \textbf{dle chuti}\\ Sladká a pálivá paprika, pepř}
        \item \parbox[t]{\linewidth}{\raggedright \spoon \textbf{1 PL}\\ Sójové omáčky (nemusí být)}
        \item \parbox[t]{\linewidth}{\raggedright \spoon \textbf{1--2 lžičky}\\ Solamylu a 3 PL vody}
        \item \parbox[t]{\linewidth}{\raggedright \piece \textbf{2 stroužky}\\ Česneku}
        \item \parbox[t]{\linewidth}{\raggedright \spoon \textbf{1 PL}\\ Bylinek (majoránka, oregano, bazalka...)}
    \end{ingredients}
    \begin{steps}
        \item Na sádle nebo oleji osmažíme nadrobno nakrájenou cibuli.
        \item Přidáme nakrájenou slaninu a společně ještě chvíli opékáme.
        \item Přidáme rajský protlak, cukr, okořeníme a chvíli podusíme. Podle potřeby podlijeme trochou vody nebo kečupem.
        \item Zahustíme solamylem rozmíchaným v trošce studené vody a ještě chvíli dusíme.
        \item Po vypnutí přidáme rozdrcený česnek a bylinky.
    \end{steps}
\end{recipe}

\begin{recipe}[DarkRed]{Pizza}
    \begin{ingredients}
        \item \parbox[t]{\linewidth}{\raggedright \weight \textbf{0,5 kg}\\ Hladké mouky}
        \item \parbox[t]{\linewidth}{\raggedright \spoon \textbf{1--2 lžíce}\\ Olivového oleje}
        \item \parbox[t]{\linewidth}{\raggedright \spoon \textbf{1 KL}\\ Soli}
        \item \parbox[t]{\linewidth}{\raggedright \piece \textbf{1 kostka}\\ Droždí}
        \item \parbox[t]{\linewidth}{\raggedright \liquid \textbf{250 ml}\\ Vlažné vody}
        \item \parbox[t]{\linewidth}{\raggedright \indefinite \textbf{na obložení}\\ Rajský protlak, česnek, oregano, sýr, šunka, žampiony...}
    \end{ingredients}
    \begin{steps}
        \item Já na pizzu používám kynuté – nekynuté těsto (recept je jinde), použít lze ale i toto základní.
        \item Ve vlažné vodě rozdrobíme kvasnice a necháme vzejít kvásek.
        \item Kvásek smícháme s moukou, zakápneme olejem a osolíme. Propracujeme hladké těsto a vytvarujeme pizzu.
        \item Povrch potřeme lehce olejem, pak rajským protlakem smíchaným s česnekem a oreganem.
        \item Obložíme dle chuti (šunka, salám, klobása, žampiony, tuňák, rajčata, zelenina), posypeme strouhaným sýrem a kořením na pizzu.
        \item Pečeme v troubě na 200 °C, dokud není pizza pěkně křupavá.
    \end{steps}
\end{recipe}

% ==========================================
% RECEPTY – OMÁČKY A HLAVNÍ CHODY
% ==========================================

\begin{recipe}[ForestGreen]{Bramborák (s masem)}
    \begin{ingredients}
        \item \parbox[t]{\linewidth}{\raggedright \weight \textbf{1 kg}\\ Brambor}
        \item \parbox[t]{\linewidth}{\raggedright \piece \textbf{2 ks}\\ Vejce}
        \item \parbox[t]{\linewidth}{\raggedright \piece \textbf{8 stroužků}\\ Česneku}
        \item \parbox[t]{\linewidth}{\raggedright \spoon \textbf{1 PL}\\ Hladké mouky a hrst strouhanky}
        \item \parbox[t]{\linewidth}{\raggedright \indefinite \textbf{koření}\\ Sůl, pepř, majoránka}
        \item \parbox[t]{\linewidth}{\raggedright \weight \textbf{2 ks}\\ Kuřecí prsa (volitelně do těsta)}
    \end{ingredients}
    \begin{steps}
        \item Brambory nastrouháme a vymačkáme z nich přebytečnou vodu.
        \item Přidáme vejce, mouku, strouhanku, drcený česnek a koření.
        \item Pokud chceme, přidáme na malé kousky nakrájené syrové kuřecí maso.
        \item Smažíme v rozpáleném oleji dozlatova z obou stran.
    \end{steps}
\end{recipe}

\begin{recipe}[BrickRed]{Rajská omáčka}
    \begin{ingredients}
        \item \parbox[t]{\linewidth}{\raggedright \weight \textbf{1 kg}\\ Hovězího masa (zadní)}
        \item \parbox[t]{\linewidth}{\raggedright \piece \textbf{2 ks}\\ Velké cibule a 2 PL másla}
        \item \parbox[t]{\linewidth}{\raggedright \liquid \textbf{2 ks}\\ Velké rajské protlaky}
        \item \parbox[t]{\linewidth}{\raggedright \piece \textbf{15 KL}\\ Cukru krupice}
        \item \parbox[t]{\linewidth}{\raggedright \indefinite \textbf{koření}\\ Skořice (2 KL), nové koření, pepř, bobkový list}
        \item \parbox[t]{\linewidth}{\raggedright \spoon \textbf{2 PL}\\ Hladké mouky}
    \end{ingredients}
    \begin{steps}
        \item Maso uvaříme v papiňáku s kořením doměkka. Vývar schováme.
        \item Na másle osmažíme cibuli, přidáme mouku a uděláme světlou jíšku.
        \item Rozředíme vývarem, přidáme protlaky, cukr a skořici.
        \item Povaříme 10 minut a přecedíme přes síto, aby byla omáčka hladká. Podáváme s knedlíkem nebo těstovinami.
    \end{steps}
\end{recipe}

\begin{recipe}[BurntOrange]{Kuře na paprice}
    \begin{ingredients}
        \item \parbox[t]{\linewidth}{\raggedright \piece \textbf{4 ks}\\ Kuřecích prsních řízků}
        \item \parbox[t]{\linewidth}{\raggedright \piece \textbf{2 ks}\\ Cibule}
        \item \parbox[t]{\linewidth}{\raggedright \liquid \textbf{1 kelímek}\\ Sladké smetany}
        \item \parbox[t]{\linewidth}{\raggedright \spoon \textbf{2 PL}\\ Mleté sladké papriky}
        \item \parbox[t]{\linewidth}{\raggedright \spoon \textbf{2 PL}\\ Másla a trocha oleje}
        \item \parbox[t]{\linewidth}{\raggedright \spoon \textbf{1--2 PL}\\ Hladké mouky}
    \end{ingredients}
    \begin{steps}
        \item Na másle a oleji osmahneme cibuli, přidáme maso na nudličky, osolíme a orestujeme.
        \item Zasypeme paprikou (krátce, aby nezhořkla), zalijeme vodou a dusíme pod pokličkou.
        \item Když je maso měkké, vmícháme smetanu s rozmíchanou moukou.
        \item Povaříme 10 minut. Podáváme s těstovinami nebo knedlíkem.
    \end{steps}
\end{recipe}

\begin{recipe}[OliveDrab]{Přírodní plátek}
    \begin{ingredients}
        \item \parbox[t]{\linewidth}{\raggedright \piece \textbf{2 ks}\\ Cibule}
        \item \parbox[t]{\linewidth}{\raggedright \piece \textbf{4 ks}\\ Kuřecích prsních řízků (nebo vepřových kotlet)}
        \item \parbox[t]{\linewidth}{\raggedright \indefinite \textbf{dle chuti}\\ Koření (např. grilovací)}
        \item \parbox[t]{\linewidth}{\raggedright \indefinite \textbf{dle chuti}\\ Čerstvý rozmarýn (můžeme vynechat)}
        \item \parbox[t]{\linewidth}{\raggedright \weight \textbf{dle potřeby}\\ Hladká mouka}
        \item \parbox[t]{\linewidth}{\raggedright \indefinite \textbf{dle chuti}\\ Pepř, sůl, olej}
    \end{ingredients}
    \begin{steps}
        \item Kuřecí prsa nakrájíme na slabší plátky, osolíme, opepříme a posypeme kořením. Pokud je směs koření slaná, maso nesolíme.
        \item Plátky obalíme v hladké mouce a oklepeme přebytečnou mouku.
        \item V pekáči na oleji zesklovatíme na kostičky nakrájenou cibuli a plátky z obou stran zprudka opečeme.
        \item Přihodíme rozmarýn, podlijeme vodou do poloviny výšky masa, přiklopíme a vložíme do trouby rozpálené na 200 °C.
        \item Postupně podléváme, aby vznikl hustší sosík. Na poslední čtvrthodinku pekáč odklopíme a pro barvu dopečeme.
    \end{steps}
\end{recipe}

\begin{recipe}[OliveGreen]{Houbová bílá omáčka}
    \begin{ingredients}
        \item \parbox[t]{\linewidth}{\raggedright \weight \textbf{500 g}\\ Masa (hovězí nebo vepřové)}
        \item \parbox[t]{\linewidth}{\raggedright \weight \textbf{300 g}\\ Čerstvých hub (kostky)}
        \item \parbox[t]{\linewidth}{\raggedright \piece \textbf{2 ks}\\ Cibule}
        \item \parbox[t]{\linewidth}{\raggedright \liquid \textbf{1 kelímek}\\ Šlehačky}
        \item \parbox[t]{\linewidth}{\raggedright \spoon \textbf{2 PL}\\ Hladké mouky}
        \item \parbox[t]{\linewidth}{\raggedright \indefinite \textbf{koření}\\ Sůl, pepř, kmín}
    \end{ingredients}
    \begin{steps}
        \item Na sádle osmažíme cibuli dozlatova, přidáme maso na kostky a orestujeme.
        \item Osolíme, okmínujeme a dusíme (u hovězího podlijeme a dusíme hodinu, pak přidáme houby; u vepřového dáváme houby hned).
        \item Do měkkého masa vlijeme smetanu s moukou.
        \item Povaříme a dochutíme. Podáváme s houskovým knedlíkem.
    \end{steps}
\end{recipe}

% ==========================================
% RECEPTY – SPECIALITY OD PŮLDY A DALŠÍ
% ==========================================

\begin{recipe}[ForestGreen]{Čína od půldy}
    \begin{ingredients}
        \item \parbox[t]{\linewidth}{\raggedright \weight \textbf{500 g}\\ Kuřecích prsou}
        \item \parbox[t]{\linewidth}{\raggedright \piece \textbf{2 ks}\\ Velké cibule}
        \item \parbox[t]{\linewidth}{\raggedright \liquid \textbf{100 ml}\\ Oleje (na naložení)}
        \item \parbox[t]{\linewidth}{\raggedright \piece \textbf{1 bal.}\\ Mražená asijská zelenina}
        \item \parbox[t]{\linewidth}{\raggedright \spoon \textbf{50 ml}\\ Ústřicové omáčky}
        \item \parbox[t]{\linewidth}{\raggedright \spoon \textbf{1 PL}\\ Sójové omáčky (nesolené)}
    \end{ingredients}
    \begin{steps}
        \item Maso nakrájíme na nudličky, obalíme v koření na čínu, zalijeme olejem a necháme uležet (nejlépe do druhého dne).
        \item Na pánvi osmahneme cibuli, přidáme naložené maso a restujeme doměkka.
        \item Přidáme zeleninu, promícháme a 5 minut dusíme pod pokličkou.
        \item Dochutíme ústřicovou a sójovou omáčkou, chilli a pepřem. Podáváme s čínskými nudlemi.
    \end{steps}
\end{recipe}

\begin{recipe}[NavyBlue]{Těstoviny s pršutem od půldy}
    \begin{ingredients}
        \item \parbox[t]{\linewidth}{\raggedright \weight \textbf{500 g}\\ Těstovin (tagliatelle)}
        \item \parbox[t]{\linewidth}{\raggedright \weight \textbf{200 g}\\ Libového pršutu}
        \item \parbox[t]{\linewidth}{\raggedright \piece \textbf{2 malé}\\ Konzervy jemného hrášku}
        \item \parbox[t]{\linewidth}{\raggedright \liquid \textbf{100 ml}\\ Olivového oleje}
        \item \parbox[t]{\linewidth}{\raggedright \weight \textbf{100 g}\\ Parmazánu (strouhaný)}
    \end{ingredients}
    \begin{steps}
        \item Těstoviny uvaříme v osolené vodě dle návodu a scedíme.
        \item Promícháme s olivovým olejem, hráškem a pršutem nakrájeným na tenké proužky.
        \item Na talíři bohatě zasypeme strouhaným parmazánem.
    \end{steps}
\end{recipe}

\begin{recipe}[RawSienna]{Halušky s brynzou od půldy}
    \begin{ingredients}
        \item \parbox[t]{\linewidth}{\raggedright \piece \textbf{2 sáčky}\\ Halušek v prášku}
        \item \parbox[t]{\linewidth}{\raggedright \weight \textbf{1 ks}\\ Brynza (ve staniolu)}
        \item \parbox[t]{\linewidth}{\raggedright \piece \textbf{2 nožičky}\\ Pikantní klobásy}
        \item \parbox[t]{\linewidth}{\raggedright \weight \textbf{300 g}\\ Slaniny a 2 cibule}
        \item \parbox[t]{\linewidth}{\raggedright \spoon \textbf{2 PL}\\ Sádla}
    \end{ingredients}
    \begin{steps}
        \item Připravíme halušky podle návodu, scedíme a propláchneme studenou vodou.
        \item Ještě horké halušky promícháme s brynzou.
        \item Na sádle osmažíme cibuli se slaninou a ke konci přidáme opéct i klobásu na půlkolečka.
        \item Halušky podáváme posypané touto smaženou směsí.
    \end{steps}
\end{recipe}

\begin{recipe}[BrickRed]{Pštrosí vejce}
    \begin{ingredients}
        \item \parbox[t]{\linewidth}{\raggedright \piece \textbf{4 ks}\\ Vajec (natvrdo)}
        \item \parbox[t]{\linewidth}{\raggedright \weight \textbf{500 g}\\ Mletého masa (mix)}
        \item \parbox[t]{\linewidth}{\raggedright \piece \textbf{2 ks}\\ Syrová vejce (do masa)}
        \item \parbox[t]{\linewidth}{\raggedright \indefinite \textbf{koření}\\ Sůl, pepř, majoránka, paprika}
        \item \parbox[t]{\linewidth}{\raggedright \spoon \textbf{2 PL}\\ Strouhanky a panák vody}
    \end{ingredients}
    \begin{steps}
        \item Mleté maso smícháme s kořením, syrovými vejci, strouhankou a vodou. Musí být lepivé.
        \item Oloupaná vařená vejce obalíme masovou směsí tak, aby vznikla koule.
        \item Dáme do vymazaného pekáčku a pečeme cca 40 minut na 180 °C.
        \item Podáváme s bramborovou kaší nebo zastudena s chlebem a hořčicí.
    \end{steps}
\end{recipe}

\begin{recipe}[Goldenrod]{Šunkofleky}
    \begin{ingredients}
        \item \parbox[t]{\linewidth}{\raggedright \weight \textbf{250 g}\\ Těstovin (fleky)}
        \item \parbox[t]{\linewidth}{\raggedright \weight \textbf{250 g}\\ Uzeného masa}
        \item \parbox[t]{\linewidth}{\raggedright \piece \textbf{4 ks}\\ Vejce}
        \item \parbox[t]{\linewidth}{\raggedright \liquid \textbf{200 ml}\\ Mléka nebo smetany}
        \item \parbox[t]{\linewidth}{\raggedright \indefinite \textbf{koření}\\ Sůl, pepř, kari}
    \end{ingredients}
    \begin{steps}
        \item Těstoviny uvaříme a dáme do vymazaného pekáče. Přidáme nakrájené uzené a koření.
        \item Dáme na 10--15 minut zapéct do trouby na 200 °C.
        \item Zalijeme vejci rozšlehanými v osoleném mléce.
        \item Pečeme, dokud vejce neztuhnou a vršek není dozlatova.
    \end{steps}
\end{recipe}

\begin{recipe}[NavyBlue]{Francouzské brambory}
    \begin{ingredients}
        \item \parbox[t]{\linewidth}{\raggedright \piece \textbf{10 ks}\\ Brambor (syrové plátky)}
        \item \parbox[t]{\linewidth}{\raggedright \weight \textbf{500 g}\\ Uzeného masa nebo točeného salámu}
        \item \parbox[t]{\linewidth}{\raggedright \piece \textbf{3 ks}\\ Cibule}
        \item \parbox[t]{\linewidth}{\raggedright \piece \textbf{3 ks}\\ Vejce}
        \item \parbox[t]{\linewidth}{\raggedright \liquid \textbf{150 ml}\\ Mléka nebo smetany}
        \item \parbox[t]{\linewidth}{\raggedright \piece \textbf{5 plátků}\\ Anglické slaniny (navrch)}
    \end{ingredients}
    \begin{steps}
        \item Do vymazaného pekáče vrstvíme: brambory (osolit/opepřit), cibuli a maso. Končíme vrstvou brambor.
        \item Pečeme v troubě na 180 °C asi 30 minut.
        \item Poté zalijeme vejci rozšlehanými v mléce a poklademe slaninou.
        \item Dopékáme další půlhodinu, dokud nejsou brambory měkké.
    \end{steps}
\end{recipe}

\begin{recipe}[Crimson]{Vinná klobása}
    \begin{ingredients}
        \item \parbox[t]{\linewidth}{\raggedright \weight \textbf{500 g}\\ Vinné klobásy}
        \item \parbox[t]{\linewidth}{\raggedright \spoon \textbf{2 PL}\\ Hladké mouky}
        \item \parbox[t]{\linewidth}{\raggedright \indefinite \textbf{koření}\\ Grilovací koření}
        \item \parbox[t]{\linewidth}{\raggedright \liquid \textbf{na smažení}\\ Olej}
    \end{ingredients}
    \begin{steps}
        \item Klobásu nakrájíme na kousky (žížalky) a obalíme v mouce smíchané s kořením.
        \item Smažíme dozlatova na pánvi.
        \item **Alternativa:** Zatočíme do šneka, propíchneme špejlí a pečeme v troubě na pečicím papíru.
    \end{steps}
\end{recipe}

% ==========================================
% RECEPTY – ZÁVĚREČNÉ SPECIALITY
% ==========================================

\begin{recipe}[ForestGreen]{Červencová smaženice z chaty \uv{Na červenci}}
    \begin{ingredients}
        \item \parbox[t]{\linewidth}{\raggedright \weight \textbf{500 g}\\ Lesních hub (směs)}
        \item \parbox[t]{\linewidth}{\raggedright \weight \textbf{250 g}\\ Anglické slaniny}
        \item \parbox[t]{\linewidth}{\raggedright \piece \textbf{1 ks}\\ Velká cibule}
        \item \parbox[t]{\linewidth}{\raggedright \piece \textbf{2 stroužky}\\ Česneku}
        \item \parbox[t]{\linewidth}{\raggedright \spoon \textbf{2 PL}\\ Sádla}
        \item \parbox[t]{\linewidth}{\raggedright \indefinite \textbf{2 KL}\\ Mletého kmínu, sůl, pepř}
    \end{ingredients}
    \begin{steps}
        \item Houby očistíme a nakrájíme na kostky. Cibuli nasekáme nadrobno.
        \item Na sádle opečeme cibuli se slaninou. Přidáme česnek a houby.
        \item Osolíme, okmínujeme a dusíme doměkka (podle potřeby lehce podlijeme).
        \item Podáváme s čerstvým chlebem.
    \end{steps}
\end{recipe}

\begin{recipe}[DarkCyan]{Špízy}
    \begin{ingredients}
        \item \parbox[t]{\linewidth}{\raggedright \weight \textbf{500 g}\\ Kuřecích prsou}
        \item \parbox[t]{\linewidth}{\raggedright \piece \textbf{2 ks}\\ Papriky (červená, žlutá)}
        \item \parbox[t]{\linewidth}{\raggedright \piece \textbf{2 nožičky}\\ Pálivé klobásy}
        \item \parbox[t]{\linewidth}{\raggedright \piece \textbf{3 ks}\\ Cibule}
        \item \parbox[t]{\linewidth}{\raggedright \indefinite \textbf{dle chuti}\\ Grilovací koření, olej}
    \end{ingredients}
    \begin{steps}
        \item Maso nakrájíme na špalíčky a naložíme přes noc do oleje s kořením.
        \item Na špejli střídavě napichujeme maso, cibuli, klobásu a papriku.
        \item Opékáme na grilu, pánvi nebo pečeme v troubě.
    \end{steps}
\end{recipe}

\begin{recipe}[Sienna]{Živáňská}
    \begin{ingredients}
        \item \parbox[t]{\linewidth}{\raggedright \weight \textbf{400 g}\\ Masa (kuřecí nebo vepřové)}
        \item \parbox[t]{\linewidth}{\raggedright \piece \textbf{4 ks}\\ Brambory (na plátky)}
        \item \parbox[t]{\linewidth}{\raggedright \piece \textbf{1 ks}\\ Klobása a 1 ks paprika}
        \item \parbox[t]{\linewidth}{\raggedright \weight \textbf{40 dkg}\\ Anglické slaniny a 1 cibule}
        \item \parbox[t]{\linewidth}{\raggedright \indefinite \textbf{dle chuti}\\ Sůl, pepř, kmín}
    \end{ingredients}
    \begin{steps}
        \item Připravíme si čtverce alobalu (raději dvojitě). Střed potřeme máslem.
        \item Vrstvíme: brambory, maso, cibuli, papriku, klobásu a navrch slaninu.
        \item Alobal pečlivě zabalíme do balíčků, aby nevytekla šťáva.
        \item Pečeme na plechu cca 1 hodinu při 200 °C.
    \end{steps}
\end{recipe}

\begin{recipe}[MidnightBlue]{Pivní placičky (Klatovské medailonky)}
    \begin{ingredients}
        \item \parbox[t]{\linewidth}{\raggedright \piece \textbf{4 ks}\\ Kuřecích řízků}
        \item \parbox[t]{\linewidth}{\raggedright \piece \textbf{2 ks}\\ Vejce (těstíčko)}
        \item \parbox[t]{\linewidth}{\raggedright \liquid \textbf{2 dcl}\\ Piva (těstíčko)}
        \item \parbox[t]{\linewidth}{\raggedright \spoon \textbf{2 PL}\\ Sójové omáčky a 2 PL solamylu}
        \item \parbox[t]{\linewidth}{\raggedright \spoon \textbf{4 PL}\\ Hladké mouky}
    \end{ingredients}
    \begin{steps}
        \item Kuřecí maso nakrájíme na malé kostičky.
        \item Smícháme suroviny na těstíčko, vložíme do něj maso a necháme přes noc uležet.
        \item Na pánvi smažíme lžící tvořené placičky dozlatova.
    \end{steps}
\end{recipe}

\begin{recipe}[Goldenrod]{Plnění šneci}
    \begin{ingredients}
        \item \parbox[t]{\linewidth}{\raggedright \weight \textbf{300 g}\\ Hladké mouky a 200 g tuku}
        \item \parbox[t]{\linewidth}{\raggedright \spoon \textbf{3 PL}\\ Zakysané smetany (těsto)}
        \item \parbox[t]{\linewidth}{\raggedright \weight \textbf{200 g}\\ Šunky (mleté/sekané)}
        \item \parbox[t]{\linewidth}{\raggedright \weight \textbf{300 g}\\ Tvrdého sýra (strouhaný)}
        \item \parbox[t]{\linewidth}{\raggedright \liquid \textbf{1/2 hrnku}\\ Zakysané smetany (náplň)}
    \end{ingredients}
    \begin{steps}
        \item Vypracujeme těsto, necháme 30 min v lednici.
        \item Náplň: smícháme sýr, šunku, vejce, pepř a smetanu.
        \item Těsto vyválíme na plát, potřeme směsí a srolujeme jako roládu.
        \item Krájíme na kolečka, klademe na plech a pečeme dozlatova.
    \end{steps}
\end{recipe}

% ==========================================
% RECEPTY – VÁNOČNÍ CUKROVÍ
% ==========================================

\begin{recipe}[BrickRed]{Išelské dortíčky od babči Libči}
    \begin{ingredients}
        \item \parbox[t]{\linewidth}{\raggedright \weight \textbf{400 g}\\ Hladké mouky}
        \item \parbox[t]{\linewidth}{\raggedright \weight \textbf{300 g}\\ Másla / Hery}
        \item \parbox[t]{\linewidth}{\raggedright \weight \textbf{140 g}\\ Moučkového cukru}
        \item \parbox[t]{\linewidth}{\raggedright \weight \textbf{60 g}\\ Mletých ořechů}
        \item \parbox[t]{\linewidth}{\raggedright \piece \textbf{2 ks}\\ Žloutky a 40 g kakaa}
        \item \parbox[t]{\linewidth}{\raggedright \indefinite \textbf{náplň}\\ Máslo, cukr, ořechy, rum}
    \end{ingredients}
    \begin{steps}
        \item Vypracujeme těsto a necháme do druhého dne v lednici uležet.
        \item Vyválíme placku (2 mm), vykrájíme kolečka a upečená (cca 150--170 °C) slepujeme náplní.
        \item Vršek ozdobíme čokoládovou polevou a půlkou mandle nebo ořechu.
    \end{steps}
\end{recipe}

\begin{recipe}[Goldenrod]{Linecké koláčky od bábí Prosecké}
    \begin{ingredients}
        \item \parbox[t]{\linewidth}{\raggedright \weight \textbf{250 g}\\ Hladké mouky}
        \item \parbox[t]{\linewidth}{\raggedright \weight \textbf{150 g}\\ Másla}
        \item \parbox[t]{\linewidth}{\raggedright \weight \textbf{100 g}\\ Moučkového cukru}
        \item \parbox[t]{\linewidth}{\raggedright \piece \textbf{1 ks}\\ Žloutek a vanilkový cukr}
        \item \parbox[t]{\linewidth}{\raggedright \indefinite \textbf{náplň}\\ Marmeláda na slepování}
    \end{ingredients}
    \begin{steps}
        \item Vypracujeme těsto, necháme uležet v chladu.
        \item Vykrajujeme plné tvary a stejný počet s vykrojeným středem.
        \item Pečeme dorůžova (cca 170 °C). Slepujeme marmeládou a sypeme cukrem.
    \end{steps}
\end{recipe}

\begin{recipe}[Fuchsia]{Citrónové cukroví od babči Libči}
    \begin{ingredients}
        \item \parbox[t]{\linewidth}{\raggedright \weight \textbf{250 g}\\ Hladké mouky}
        \item \parbox[t]{\linewidth}{\raggedright \weight \textbf{150 g}\\ Másla a 150 g cukru}
        \item \parbox[t]{\linewidth}{\raggedright \piece \textbf{1 ks}\\ Žloutek}
        \item \parbox[t]{\linewidth}{\raggedright \liquid \textbf{1 ks}\\ Šťáva a kůra z citronu}
        \item \parbox[t]{\linewidth}{\raggedright \indefinite \textbf{poleva}\\ Cukr, citron. šťáva, voda}
    \end{ingredients}
    \begin{steps}
        \item Vypracujeme těsto, necháme uležet, vykrájíme a upečeme dorůžova.
        \item Slepujeme marmeládou.
        \item Potřeme citronovou polevou (cukr zašlehaný se šťávou a horkou vodou).
    \end{steps}
\end{recipe}

\begin{recipe}[Sienna]{Ořechy od babči Libči}
    \begin{ingredients}
        \item \parbox[t]{\linewidth}{\raggedright \weight \textbf{300 g}\\ Hladké mouky}
        \item \parbox[t]{\linewidth}{\raggedright \weight \textbf{200 g}\\ Másla}
        \item \parbox[t]{\linewidth}{\raggedright \weight \textbf{150 g}\\ Moučkového cukru}
        \item \parbox[t]{\linewidth}{\raggedright \weight \textbf{160 g}\\ Mletých ořechů}
        \item \parbox[t]{\linewidth}{\raggedright \spoon \textbf{2 PL}\\ Kakaa a 2 KL skořice}
        \item \parbox[t]{\linewidth}{\raggedright \piece \textbf{1 ks}\\ Vejce}
    \end{ingredients}
    \begin{steps}
        \item Vypracujeme těsto a necháme uležet. Formičky vymažeme tukem.
        \item Těsto vtlačíme do formiček tak, aby uprostřed zbyl důlek na náplň.
        \item Upečeme, vyklopíme a plníme ořechovým krémem (stejný jako u Išelských).
    \end{steps}
\end{recipe}

\begin{recipe}[NavyBlue]{Rohlíčky od tety Aleny}
    \begin{ingredients}
        \item \parbox[t]{\linewidth}{\raggedright \weight \textbf{300 g}\\ Hladké mouky}
        \item \parbox[t]{\linewidth}{\raggedright \weight \textbf{200 g}\\ Másla}
        \item \parbox[t]{\linewidth}{\raggedright \weight \textbf{100 g}\\ Mletých ořechů}
        \item \parbox[t]{\linewidth}{\raggedright \weight \textbf{40 g}\\ Moučkového cukru}
        \item \parbox[t]{\linewidth}{\raggedright \piece \textbf{2 ks}\\ Žloutky a vanilkový cukr}
    \end{ingredients}
    \begin{steps}
        \item Zpracujeme těsto, necháme v chladu odležet.
        \item Tvarujeme malé rohlíčky a pečeme na plechu dorůžova.
        \item Ještě horké je opatrně obalujeme v moučkovém cukru smíchaném s vanilkovým.
    \end{steps}
\end{recipe}

\begin{recipe}[SaddleBrown]{Medvědí tlapičky od babči Libči}
    \begin{ingredients}
        \item \parbox[t]{\linewidth}{\raggedright \weight \textbf{200 g}\\ Másla}
        \item \parbox[t]{\linewidth}{\raggedright \weight \textbf{200 g}\\ Ořechů}
        \item \parbox[t]{\linewidth}{\raggedright \weight \textbf{200 g}\\ Hladké mouky}
        \item \parbox[t]{\linewidth}{\raggedright \weight \textbf{100 g}\\ Moučkového cukru}
        \item \parbox[t]{\linewidth}{\raggedright \weight \textbf{60 g}\\ Strouhané čokolády na vaření}
        \item \parbox[t]{\linewidth}{\raggedright \weight \textbf{1/2 KL}\\ Tlučeného hřebíčku}
        \item \parbox[t]{\linewidth}{\raggedright \piece \textbf{1 ks}\\ Celé vejce}
        \item \parbox[t]{\linewidth}{\raggedright \weight \textbf{50 g}\\ Hery na vymazání formiček}
        \item \parbox[t]{\linewidth}{\raggedright \indefinite \textbf{dle potřeby}\\ Moučkový a vanilkový cukr na obalení}
    \end{ingredients}
    \begin{steps}
        \item Ze všech surovin vypracujeme těsto, zabalíme do folie a necháme nejlépe do druhého dne uležet.
        \item Rozehřejeme heru a vymažeme formičky. Těsto do nich vmáčkneme tak, aby uprostřed byla prohlubeň.
        \item Upečeme, vyklopíme a ještě horké obalujeme v moučkovém cukru smíchaném s vanilkovým.
    \end{steps}
\end{recipe}

\begin{recipe}[LightSkyBlue]{Sněhová kolečka od maminky}
    \begin{ingredients}
        \item \parbox[t]{\linewidth}{\raggedright \piece \textbf{4 ks}\\ Bílky}
        \item \parbox[t]{\linewidth}{\raggedright \weight \textbf{320 g}\\ Moučkového cukru}
        \item \parbox[t]{\linewidth}{\raggedright \spoon \textbf{1 PL}\\ Citronové šťávy}
        \item \parbox[t]{\linewidth}{\raggedright \indefinite \textbf{dle potřeby}\\ Čokoládová poleva}
    \end{ingredients}
    \begin{steps}
        \item V míse vyšleháme elektrickým šlehačem na nejvyšším stupni tuhý sníh z bílků, cukru a šťávy.
        \item Cukrářským sáčkem tvoříme na plechu vyloženém pečícím papírem věnečky.
        \item Sušíme v předehřáté troubě cca 1 hodinu (el. trouba 90 °C, horkovzduch 70 °C).
        \item Po vychladnutí můžeme věnečky máčet do čokoládové polevy.
    \end{steps}
\end{recipe}

\begin{recipe}[HotPink]{Mušličky od Simony Tobiškové starší}
    \begin{ingredients}
        \item \parbox[t]{\linewidth}{\raggedright \weight \textbf{280 g}\\ Hladké mouky}
        \item \parbox[t]{\linewidth}{\raggedright \weight \textbf{280 g}\\ Másla}
        \item \parbox[t]{\linewidth}{\raggedright \piece \textbf{2 ks}\\ Celá vejce}
        \item \parbox[t]{\linewidth}{\raggedright \weight \textbf{110 g}\\ Oříšků (náplň)}
        \item \parbox[t]{\linewidth}{\raggedright \weight \textbf{110 g}\\ Moučkového cukru (náplň)}
        \item \parbox[t]{\linewidth}{\raggedright \piece \textbf{1 ks}\\ Celé vejce a 1 žloutek (náplň)}
        \item \parbox[t]{\linewidth}{\raggedright \indefinite \textbf{dle potřeby}\\ Moučkový cukr na obalení}
    \end{ingredients}
    \begin{steps}
        \item Ze surovin vypracujeme těsto, zabalíme do folie a necháme nejlépe do druhého dne uležet.
        \item Suroviny na náplň v misce utřeme.
        \item Těsto vyválíme a vykrajujeme kolečka, dáme na ně náplň, přehneme na polovinu a přimáčkneme.
        \item Upečeme a ještě horké obalujeme v moučkovém cukru.
    \end{steps}
\end{recipe}

\begin{recipe}[Teal]{Myslivecké knoflíky od Simony Tobiškové starší}
    \begin{ingredients}
        \item \parbox[t]{\linewidth}{\raggedright \weight \textbf{300 g}\\ Hladké mouky}
        \item \parbox[t]{\linewidth}{\raggedright \weight \textbf{140 g}\\ Másla}
        \item \parbox[t]{\linewidth}{\raggedright \weight \textbf{40 g}\\ Moučkového cukru}
        \item \parbox[t]{\linewidth}{\raggedright \piece \textbf{3 ks}\\ Žloutky}
        \item \parbox[t]{\linewidth}{\raggedright \piece \textbf{3 ks}\\ Bílky (náplň)}
        \item \parbox[t]{\linewidth}{\raggedright \weight \textbf{110 g}\\ Moučkového cukru (náplň)}
        \item \parbox[t]{\linewidth}{\raggedright \weight \textbf{110 g}\\ Oříšků (náplň)}
        \item \parbox[t]{\linewidth}{\raggedright \indefinite \textbf{dle potřeby}\\ Půlky vlašských ořechů na vmáčknutí}
    \end{ingredients}
    \begin{steps}
        \item Ze surovin vypracujeme těsto, zabalíme do folie a necháme nejlépe do druhého dne uležet.
        \item Těsto vyválíme a vykrajujeme kolečka.
        \item Na každé kolečko dáme doprostřed trochu náplně (do ušlehaných bílků lehce vmícháme cukr a ořechy) a zmáčkneme polovinou ořechu.
        \item Připravené knoflíky upečeme.
    \end{steps}
\end{recipe}

\begin{recipe}[DarkOrchid]{Nuteláčci od maminky}
    \begin{ingredients}
        \item \parbox[t]{\linewidth}{\raggedright \weight \textbf{250 g}\\ Hladké mouky}
        \item \parbox[t]{\linewidth}{\raggedright \weight \textbf{160 g}\\ Másla a 100 g cukru}
        \item \parbox[t]{\linewidth}{\raggedright \spoon \textbf{2 PL}\\ Kakaa a 1 žloutek}
        \item \parbox[t]{\linewidth}{\raggedright \indefinite \textbf{na ozdobu}\\ Nutela, sekané ořechy}
        \item \parbox[t]{\linewidth}{\raggedright \indefinite \textbf{extra}\\ Čokoláda a minilentilky}
    \end{ingredients}
    \begin{steps}
        \item Upečeme linecká kolečka z kakaového těsta.
        \item Slepíme dvě kolečka k sobě nutelou.
        \item Horní stranu mázneme nutelou, vmáčkneme do ořechů, namočíme do čokolády a ozdobíme lentilkou.
    \end{steps}
\end{recipe}

\begin{recipe}[Maroon]{Vosí hnízda}
    \begin{ingredients}
        \item \parbox[t]{\linewidth}{\raggedright \weight \textbf{200 g}\\ Piškotů (drcených)}
        \item \parbox[t]{\linewidth}{\raggedright \weight \textbf{100 g}\\ Másla a 100 g oříšků}
        \item \parbox[t]{\linewidth}{\raggedright \weight \textbf{180 g}\\ Moučkového cukru}
        \item \parbox[t]{\linewidth}{\raggedright \spoon \textbf{3 PL}\\ Kakaa a 4 PL rumu}
        \item \parbox[t]{\linewidth}{\raggedright \indefinite \textbf{náplň}\\ Máslo, cukr, žloutek/koňak}
        \item \parbox[t]{\linewidth}{\raggedright \piece \textbf{balení}\\ Kulaté piškoty (základna)}
    \end{ingredients}
    \begin{steps}
        \item Smícháme drcené piškoty se surovinami na těsto. Formičky vysypeme cukrem.
        \item Vtlačíme těsto, uděláme vařečkou důlek, vyklepneme.
        \item Do důlku dáme krém a přitiskneme piškot. Špičky můžeme máčet v čokoládě.
    \end{steps}
\end{recipe}

\begin{recipe}[DodgerBlue]{Kokosové protlačované cukroví}
    \begin{ingredients}
        \item \parbox[t]{\linewidth}{\raggedright \weight \textbf{380 g}\\ Hladké mouky}
        \item \parbox[t]{\linewidth}{\raggedright \weight \textbf{110 g}\\ Moučkového cukru}
        \item \parbox[t]{\linewidth}{\raggedright \piece \textbf{2 ks}\\ Celá vejce}
        \item \parbox[t]{\linewidth}{\raggedright \weight \textbf{100 g}\\ Mletého kokosu}
        \item \parbox[t]{\linewidth}{\raggedright \weight \textbf{210 g}\\ Másla}
        \item \parbox[t]{\linewidth}{\raggedright \weight \textbf{100 g}\\ Čokolády na vaření (poleva)}
        \item \parbox[t]{\linewidth}{\raggedright \weight \textbf{2 cm}\\ Omegy nebo jiného 100\% tuku (poleva)}
    \end{ingredients}
    \begin{steps}
        \item Na vále vypracujeme těsto, zabalíme do folie a necháme nejlépe do druhého dne uležet.
        \item Těsto protlačujeme na masovém strojku přes speciální nástavec a tvoříme kolečka, esíčka, účka nebo tyčinky.
        \item Upečeme do světle růžova.
        \item Konce tvarů máčíme v čokoládové polevě.
    \end{steps}
\end{recipe}

\begin{recipe}[PowderBlue]{Kokosky (obyčejné)}
    \begin{ingredients}
        \item \parbox[t]{\linewidth}{\raggedright \piece \textbf{5 ks}\\ Bílků}
        \item \parbox[t]{\linewidth}{\raggedright \piece \textbf{1 balíček}\\ Mletého kokosu}
        \item \parbox[t]{\linewidth}{\raggedright \weight \textbf{100 g}\\ Moučkového cukru}
    \end{ingredients}
    \begin{steps}
        \item Z bílků ušleháme tuhý sníh, do kterého zašleháme cukr.
        \item Lehce vmícháme kokos.
        \item Na plech děláme lžičkou kopečky a opatrně upečeme.
    \end{steps}
\end{recipe}

\begin{recipe}[PaleTurquoise]{Kokosky luxusní s delší přípravou (recept z internetu)}
    \begin{ingredients}
        \item \parbox[t]{\linewidth}{\raggedright \piece \textbf{3 ks}\\ Bílků}
        \item \parbox[t]{\linewidth}{\raggedright \weight \textbf{180 g}\\ Moučkového cukru}
        \item \parbox[t]{\linewidth}{\raggedright \weight \textbf{150--180 g}\\ Strouhaného kokosu (dle kvality kokosu)}
    \end{ingredients}
    \begin{steps}
        \item Z bílků ušleháme měkký, řídký sníh.
        \item Misku se sněhem postavíme do vodní lázně nad páru (voda těsně pod bodem varu, nesmí vřít). Postupně zašleháme cukr.
        \item Šleháme nad párou, až je hmota příjemně horká. Bílky nenecháme srazit, stále mícháme a šleháme (3–5 minut).
        \item Uvařený sníh odstavíme a necháme zchladnout, občas prošleháme.
        \item Do vychladlého sněhu jemně zašleháme kokos. Nesmí ho být příliš, hmota by se začala drolit a byla příliš suchá.
        \item Troubu předehřejeme na 140 °C, plechy vyložíme papírem. Lžičkou vykrajujeme malé hromádky a pečeme 15--20 minut do jemného zrůžovění. Necháme zchladnout na plechu.
        \item Vychladlé kokosky necháme pár dní rozležet, cosi se na nich změní k lepšímu.
    \end{steps}
\end{recipe}

\begin{recipe}[OrangeRed]{Rychlé koblížky (babičky Jižňácké)}
    \begin{ingredients}
        \item \parbox[t]{\linewidth}{\raggedright \weight \textbf{250 g}\\ Polohrubé mouky}
        \item \parbox[t]{\linewidth}{\raggedright \piece \textbf{1 ks}\\ Tvaroh (vanička)}
        \item \parbox[t]{\linewidth}{\raggedright \piece \textbf{3 ks}\\ Vejce a 1 prdopeč}
        \item \parbox[t]{\linewidth}{\raggedright \spoon \textbf{3 PL}\\ Moučkového cukru}
        \item \parbox[t]{\linewidth}{\raggedright \liquid \textbf{5 PL}\\ Mléka a 5 PL rumu}
    \end{ingredients}
    \begin{steps}
        \item Všechny suroviny smícháme v misce.
        \item Lžící děláme hromádky přímo do rozpáleného oleje.
        \item Smažíme dorůžova. Obalujeme v cukru nebo zdobíme marmeládou.
    \end{steps}
\end{recipe}

\begin{recipe}[BurntOrange]{Marokánky}
    \begin{ingredients}
        \item \parbox[t]{\linewidth}{\raggedright \liquid \textbf{0,5 l}\\ Mléka a 100 g másla}
        \item \parbox[t]{\linewidth}{\raggedright \weight \textbf{8 PL}\\ Hladké mouky a 10 PL cukru}
        \item \parbox[t]{\linewidth}{\raggedright \indefinite \textbf{směs}\\ Rozinky, mandle, ořechy}
        \item \parbox[t]{\linewidth}{\raggedright \piece \textbf{100 g}\\ Kandované ovoce}
        \item \parbox[t]{\linewidth}{\raggedright \indefinite \textbf{poleva}\\ Čokoláda na vaření}
    \end{ingredients}
    \begin{steps}
        \item Mléko, mouku a cukr vaříme do zhoustnutí (kaše), pak přidáme máslo.
        \item Do hmoty vmícháme nasekané ořechy, mandle a ovoce.
        \item Na plech děláme lžičkou placičky a pečeme na 180 °C, až okraje zhnědnou.
        \item Po vychladnutí spodní rovnou stranu potíráme čokoládou.
    \end{steps}
\end{recipe}

\begin{recipe}[Pink]{Ovocné košíčky}
    \begin{ingredients}
        \item \parbox[t]{\linewidth}{\raggedright \weight \textbf{250 g}\\ Hladké mouky (těsto)}
        \item \parbox[t]{\linewidth}{\raggedright \weight \textbf{150 g}\\ Tuku (těsto)}
        \item \parbox[t]{\linewidth}{\raggedright \weight \textbf{120 g}\\ Moučkového cukru (těsto)}
        \item \parbox[t]{\linewidth}{\raggedright \piece \textbf{2 ks}\\ Žloutky (těsto)}
        \item \parbox[t]{\linewidth}{\raggedright \weight \textbf{1/2 KL}\\ Prášku do pečiva (těsto)}
        \item \parbox[t]{\linewidth}{\raggedright \piece \textbf{1 balíček}\\ Vanilkového pudinku (krém)}
        \item \parbox[t]{\linewidth}{\raggedright \liquid \textbf{300 ml}\\ Mléka (krém)}
        \item \parbox[t]{\linewidth}{\raggedright \weight \textbf{80 g}\\ Másla a 80 g moučkového cukru (krém)}
        \item \parbox[t]{\linewidth}{\raggedright \spoon \textbf{2 PL}\\ Rumu (krém)}
        \item \parbox[t]{\linewidth}{\raggedright \indefinite \textbf{dle sezóny}\\ Ovoce (borůvky, jahody, kiwi...)}
    \end{ingredients}
    \begin{steps}
        \item Mouku, cukr a prášek do pečiva dobře promícháme, přidáme změklý tuk, žloutky a vypracujeme vláčné těsto. Dáme na půl hodiny do lednice.
        \item Formičky vymažeme tukem a naplníme těstem v ne příliš silné vrstvě (těsto pečením nabude). Upečeme ve středně vyhřáté troubě.
        \item Z mléka a pudinkového prášku uvaříme hustý pudink. Do ještě vlažného pudinku zašleháme máslo a cukr, přidáme rum.
        \item Košíčky naplníme krémem a ozdobíme ovocem.
    \end{steps}
\end{recipe}

\begin{recipe}[MediumOrchid]{Ořechové trubičky}
    \begin{ingredients}
        \item \parbox[t]{\linewidth}{\raggedright \weight \textbf{360 g}\\ Hladké mouky (těsto)}
        \item \parbox[t]{\linewidth}{\raggedright \weight \textbf{140 g}\\ Změklého másla (těsto)}
        \item \parbox[t]{\linewidth}{\raggedright \weight \textbf{120 g}\\ Mletých lískových oříšků (těsto)}
        \item \parbox[t]{\linewidth}{\raggedright \weight \textbf{100 g}\\ Moučkového cukru (těsto)}
        \item \parbox[t]{\linewidth}{\raggedright \piece \textbf{1 ks}\\ Vejce (těsto)}
        \item \parbox[t]{\linewidth}{\raggedright \piece \textbf{1 ks}\\ Vejce a nasekané oříšky (na povrch)}
        \item \parbox[t]{\linewidth}{\raggedright \weight \textbf{200 g}\\ Čokolády na vaření (krém)}
        \item \parbox[t]{\linewidth}{\raggedright \weight \textbf{200 g}\\ Nutelly (krém)}
        \item \parbox[t]{\linewidth}{\raggedright \weight \textbf{200 g}\\ Změklého másla (krém)}
        \item \parbox[t]{\linewidth}{\raggedright \spoon \textbf{2 lžíce}\\ Mandlového nebo ořechového likéru (krém)}
    \end{ingredients}
    \begin{steps}
        \item Ze všech surovin vypracujeme těsto, které necháme v chladu odpočinout.
        \item Těsto rozválíme a rádýlkem rozdělíme na obdélníčky, které obtočíme kolem malých trubiček na kremrole.
        \item Dobře spojíme, potřeme rozšlehaným vejcem a obalíme v nahrubo nasekaných oříšcích.
        \item Pečeme při 170--180 °C dorůžova, pomocí kleští teplé stáhneme z formiček.
        \item Na krém rozpustíme čokoládu, necháme vychladnout na pokojovou teplotu a vyšleháme s máslem, likérem a nutellou na hladký krém.
        \item Trubičky naplníme a uchováváme v chladnu.
    \end{steps}
\end{recipe}

\begin{recipe}[ForestGreen]{Tortilla}
    \begin{ingredients}
        \item \parbox[t]{\linewidth}{\raggedright \piece \textbf{2 hrnky}\\ Hladké mouky}
        \item \parbox[t]{\linewidth}{\raggedright \liquid \textbf{1/2 hrnku}\\ Vlažné vody}
        \item \parbox[t]{\linewidth}{\raggedright \spoon \textbf{3 PL}\\ Oleje a 1 KL soli}
    \end{ingredients}
    \begin{steps}
        \item Vypracujeme těsto, aby nelepilo. Rozdělíme na cca 10 bochánků.
        \item Vyválíme na co nejtenčí placky.
        \item Pečeme krátce na suché rozpálené pánvi z obou stran.
    \end{steps}
\end{recipe}

\end{document}